%% ============================================================================
%%  ĐỒ ÁN TỐT NGHIỆP
%%  XÂY DỰNG VÀ ĐÁNH GIÁ HIỆU NĂNG HỆ THỐNG FOG COMPUTING CHO IOT NHÀ THÔNG MINH:
%%  SO SÁNH KIẾN TRÚC MONOLITHIC VÀ FOG TRÊN NỀN TẢNG APACHE STORM
%%
%%  Biên dịch:  pdflatex THESIS.tex  (chạy 2-3 lần để cập nhật mục lục/tham chiếu)
%%  Yêu cầu:    bản TeX có gói vntex (T5) — Overleaf đã có sẵn.
%%  Ghi chú font tiếng Việt: nếu pdflatex báo lỗi font T5, chuyển sang biên dịch
%%              bằng XeLaTeX và thay khối font bằng fontspec + Times New Roman.
%%  Ảnh Grafana đặt trong:  images/mono/  images/fog/kb1a/  images/fog/kb1b/
%%                          images/fog/kb2a/  images/fog/kb2b/
%% ============================================================================
\documentclass[13pt,a4paper]{report}

%% ---------- Mã hoá & ngôn ngữ ----------
\usepackage[T5]{fontenc}
\usepackage[utf8]{inputenc}
\usepackage[vietnamese]{babel}

%% ---------- Bố cục trang ----------
\usepackage[left=3cm,right=2.5cm,top=2.5cm,bottom=2.5cm]{geometry}
\usepackage{setspace}
\onehalfspacing

%% ---------- Đồ hoạ & bảng ----------
\usepackage{graphicx}
\usepackage{float}
\usepackage{booktabs}
\usepackage{longtable}
\usepackage{multirow}
\usepackage{array}
\usepackage{makecell}
\usepackage{xcolor}
\usepackage{caption}
\usepackage{subcaption}
\usepackage{siunitx}

%% ---------- Toán & mã nguồn ----------
\usepackage{amsmath}
\usepackage{amssymb}
\usepackage{listings}

%% ---------- Biểu đồ TikZ / pgfplots ----------
\usepackage{tikz}
\usepackage{pgfplots}
\pgfplotsset{compat=1.17}
\usepgfplotslibrary{groupplots}
\usetikzlibrary{arrows.meta,positioning,fit,backgrounds,calc,shapes.geometric,shapes.misc}
\usepackage{pgf-pie}

%% ---------- Tiêu đề chương, header/footer, mục lục ----------
\usepackage{titlesec}
\usepackage{fancyhdr}
\usepackage{tocloft}

%% ---------- Liên kết & tham chiếu (nạp gần cuối) ----------
\usepackage[hidelinks]{hyperref}
\usepackage[noabbrev,nameinlink]{cleveref}

%% cleveref chưa hỗ trợ tiếng Việt -> tự định nghĩa tên tham chiếu
\crefname{chapter}{Chương}{Chương}
\Crefname{chapter}{Chương}{Chương}
\crefname{section}{Mục}{Mục}
\Crefname{section}{Mục}{Mục}
\crefname{figure}{Hình}{Hình}
\Crefname{figure}{Hình}{Hình}
\crefname{table}{Bảng}{Bảng}
\Crefname{table}{Bảng}{Bảng}
\crefname{equation}{Công thức}{Công thức}
\Crefname{equation}{Công thức}{Công thức}

\setlength{\headheight}{15pt}

%% ============================================================================
%%  ĐỊNH NGHĨA MÀU, KIỂU MÃ NGUỒN, LỆNH TIỆN ÍCH
%% ============================================================================
\definecolor{monoRed}{RGB}{197,48,48}
\definecolor{fog8Green}{RGB}{34,139,34}
\definecolor{fog1Orange}{RGB}{221,128,18}
\definecolor{accentBlue}{RGB}{27,94,160}
\definecolor{codeBg}{RGB}{246,246,248}
\definecolor{codeKw}{RGB}{0,90,160}
\definecolor{codeCm}{RGB}{90,120,90}
\definecolor{boxGray}{RGB}{236,238,242}

\lstset{
  backgroundcolor=\color{codeBg},
  basicstyle=\ttfamily\footnotesize,
  keywordstyle=\color{codeKw}\bfseries,
  commentstyle=\color{codeCm}\itshape,
  stringstyle=\color{monoRed},
  numbers=left, numberstyle=\tiny\color{gray}, numbersep=8pt,
  breaklines=true, frame=single, rulecolor=\color{gray!40},
  showstringspaces=false, tabsize=2, captionpos=b,
  extendedchars=false, inputencoding=utf8
}

\newcommand{\code}[1]{\texttt{\small #1}}
% Ảnh Grafana: hiện ảnh nếu đã có; nếu chưa (đang đo lại) thì hiện khung placeholder
% để báo cáo VẪN biên dịch được. Sau khi đo, chỉ cần copy ảnh vào images/ là tự hiện.
\newcommand{\grafana}[1]{%
  \IfFileExists{#1}{\includegraphics[width=\linewidth]{#1}}%
  {\fbox{\parbox[c][3cm][c]{0.88\linewidth}{\centering\small
   [Ảnh Grafana sẽ chèn sau khi đo lại]\\[3pt] \texttt{\footnotesize \detokenize{#1}}}}}%
}
\sisetup{group-separator={.},group-minimum-digits=4}

%% ---------- Định dạng tiêu đề chương kiểu "CHƯƠNG 1" ----------
\titleformat{\chapter}[display]
  {\normalfont\Large\bfseries\filcenter}
  {\MakeUppercase{\chaptertitlename}\ \thechapter}{12pt}{\LARGE}
\titlespacing*{\chapter}{0pt}{10pt}{28pt}

%% ---------- Header / Footer ----------
\pagestyle{fancy}
\fancyhf{}
\renewcommand{\headrulewidth}{0.4pt}
\fancyhead[L]{\nouppercase{\leftmark}}
\fancyhead[R]{\thepage}
\fancypagestyle{plain}{\fancyhf{}\fancyhead[R]{\thepage}\renewcommand{\headrulewidth}{0pt}}

%% ---------- Tên các danh mục ----------
\addto\captionsvietnamese{%
  \renewcommand{\contentsname}{MỤC LỤC}%
  \renewcommand{\listfigurename}{DANH MỤC HÌNH ẢNH}%
  \renewcommand{\listtablename}{DANH MỤC BẢNG}%
  \renewcommand{\bibname}{TÀI LIỆU THAM KHẢO}%
  \renewcommand{\chaptername}{CHƯƠNG}%
  \renewcommand{\appendixname}{PHỤ LỤC}%
}

%% ============================================================================
\begin{document}
%% ============================================================================

%% ----------------------------------------------------------------------------
%%  TRANG BÌA
%% ----------------------------------------------------------------------------
\begin{titlepage}
\centering
\thispagestyle{empty}
{\large ĐẠI HỌC QUỐC GIA HÀ NỘI\par}
{\large TRƯỜNG ĐẠI HỌC CÔNG NGHỆ\par}
{\large KHOA CÔNG NGHỆ THÔNG TIN\par}
\vspace{1.2cm}
% Đặt logo trường (logo-uet.png) cùng thư mục để hiển thị; nếu không có thì bỏ trống.
\IfFileExists{logo-uet.png}{\includegraphics[width=0.18\textwidth]{logo-uet.png}\par}{\vspace{2.4cm}}
\vspace{0.6cm}
{\Large\bfseries ĐỒ ÁN TỐT NGHIỆP\par}
\vspace{1.6cm}
{\LARGE\bfseries XÂY DỰNG VÀ ĐÁNH GIÁ HIỆU NĂNG\\[0.3em]
HỆ THỐNG FOG COMPUTING CHO IOT NHÀ THÔNG MINH:\\[0.3em]
SO SÁNH KIẾN TRÚC MONOLITHIC VÀ FOG\\[0.3em]
TRÊN NỀN TẢNG APACHE STORM\par}
\vspace{2.2cm}
\begin{flushright}
\begin{minipage}{0.62\textwidth}
\large
\textbf{Sinh viên thực hiện:} Nguyễn Gia Bảo\\[0.4em]
\textbf{Mã sinh viên:} 22025514\\[0.4em]
\textbf{Lớp:} \dotfill\\[0.4em]
\textbf{Giảng viên hướng dẫn:}\\[0.2em]
\hspace*{1.2em}PGS.TS. Phạm Mạnh Linh\\[0.2em]
\hspace*{1.2em}ThS. Trần Anh Tú\\[0.4em]
\end{minipage}
\end{flushright}
\vfill
{\large Hà Nội, tháng 6 năm 2026\par}
\end{titlepage}

%% Chú thích: thay \includegraphics{example-image-a} bằng logo trường (logo-uit.png).

%% ----------------------------------------------------------------------------
%%  CÁC TRANG ĐẦU (đánh số La Mã)
%% ----------------------------------------------------------------------------
\pagenumbering{roman}

%% ---------- Trang xác nhận của giảng viên ----------
\chapter*{NHẬN XÉT CỦA GIẢNG VIÊN HƯỚNG DẪN}
\addcontentsline{toc}{chapter}{Nhận xét của giảng viên hướng dẫn}
\thispagestyle{plain}
\vspace{1em}
\noindent\dotfill\par\vspace{1.4em}
\noindent\dotfill\par\vspace{1.4em}
\noindent\dotfill\par\vspace{1.4em}
\noindent\dotfill\par\vspace{1.4em}
\noindent\dotfill\par\vspace{1.4em}
\noindent\dotfill\par\vspace{1.4em}
\noindent\dotfill\par\vspace{1.4em}
\noindent\dotfill\par
\vspace{2.5em}
\begin{flushright}
\textit{Hà Nội, ngày \dots\ tháng \dots\ năm 2026}\\[0.4em]
\textbf{Giảng viên hướng dẫn}\\[3.2em]
\dotfill
\end{flushright}

%% ---------- Lời cam đoan ----------
\chapter*{LỜI CAM ĐOAN}
\addcontentsline{toc}{chapter}{Lời cam đoan}
\thispagestyle{plain}
Tôi xin cam đoan đồ án tốt nghiệp với đề tài \textit{``Xây dựng và đánh giá
hiệu năng hệ thống Fog Computing cho IoT Nhà Thông Minh: So sánh kiến trúc
Monolithic và Fog trên nền tảng Apache Storm''} là công trình nghiên cứu của
riêng tôi dưới sự hướng dẫn của giảng viên hướng dẫn. Các số liệu, kết quả đo
lường trình bày trong đồ án là trung thực, được thu thập trực tiếp từ hệ thống
thực nghiệm do tôi triển khai trên nền tảng Apache Storm, Docker, Amazon EC2 và
thiết bị Raspberry Pi 3. Mọi tham khảo từ tài liệu của tác giả khác đều được
trích dẫn đầy đủ trong phần Tài liệu tham khảo. Tôi xin chịu hoàn toàn trách
nhiệm về tính trung thực của nội dung đồ án này.

\vspace{2.5em}
\begin{flushright}
\textit{Hà Nội, ngày \dots\ tháng \dots\ năm 2026}\\[0.4em]
\textbf{Sinh viên thực hiện}\\[3.2em]
Nguyễn Gia Bảo
\end{flushright}

%% ---------- Lời cảm ơn ----------
\chapter*{LỜI CẢM ƠN}
\addcontentsline{toc}{chapter}{Lời cảm ơn}
\thispagestyle{plain}
Tôi xin gửi lời cảm ơn chân thành đến quý thầy cô Khoa Công nghệ Thông tin,
Trường Đại học Công nghệ --- Đại học Quốc gia Hà Nội đã truyền đạt những
kiến thức nền tảng quý báu trong suốt quá trình học tập. Đặc biệt, tôi xin bày
tỏ lòng biết ơn sâu sắc đến giảng viên hướng dẫn đã tận tình định hướng, góp ý
và đồng hành cùng tôi trong toàn bộ quá trình thực hiện đồ án này. Tôi cũng xin
cảm ơn gia đình và bạn bè đã luôn động viên, hỗ trợ để tôi hoàn thành đồ án tốt
nghiệp.

%% ---------- Tóm tắt tiếng Việt ----------
\chapter*{TÓM TẮT}
\addcontentsline{toc}{chapter}{Tóm tắt}
\thispagestyle{plain}
Sự bùng nổ của Internet of Things (IoT) đặt ra thách thức lớn về băng thông và
độ trễ khi toàn bộ dữ liệu cảm biến được đẩy lên Cloud để xử lý. Đồ án này
nghiên cứu, thiết kế và đánh giá hiệu năng một hệ thống \textbf{Fog Computing}
cho bài toán giám sát tiêu thụ điện của 40 ngôi nhà thông minh, với tổng tải lên
tới \textbf{400 msg/s}, được xây dựng trên nền tảng Apache Storm. Hệ thống được
so sánh trực tiếp với kiến trúc \textbf{Monolithic} (xử lý toàn bộ trên Cloud)
qua sáu kịch bản thực nghiệm.

Kết quả đo lường cho thấy: kiến trúc Monolithic bão hoà tại khoảng 20 nhà
(200 msg/s), với \code{Bolt\_split} đạt \textbf{169\%} capacity và complete
latency \textbf{5,28 s} tại 400 msg/s. Ngược lại, kiến trúc Fog với
\textbf{8 gateway} duy trì cloud bolt capacity \textbf{0\%} tại 400 msg/s, gateway
xử lý ở mức sub-millisecond (\textless 0,1 ms execute latency). Cơ chế
Store-and-Forward đạt \textbf{0\% mất mát dữ liệu} khi Cloud ngắt kết nối 5 phút,
phục hồi hoàn toàn trong \textbf{71--96 s}. Lưu lượng WAN giảm khoảng
\textbf{86\%} nhờ nén GZIP theo lô. Gateway được triển khai thực tế trên thiết bị
Raspberry Pi 3. Đồ án đề xuất tổng cộng \textbf{mười cải tiến kỹ thuật} so với kiến
trúc Monolithic (cùng một bản vá lỗi đo lường mang tính tiền đề), trong đó có cơ
chế \textit{Energy Optimizer} tự động ``đóng băng'' gateway nhàn rỗi để đưa CPU về
0\%.

\vspace{0.8em}
\noindent\textbf{Từ khoá:} Fog Computing, Edge Computing, IoT, Apache Storm,
Smart Home, Store-and-Forward, Raspberry Pi, MQTT.

%% ---------- Tóm tắt tiếng Anh ----------
\chapter*{ABSTRACT}
\addcontentsline{toc}{chapter}{Abstract}
\thispagestyle{plain}
The rapid growth of the Internet of Things (IoT) raises serious bandwidth and
latency challenges when all sensor data is pushed to the Cloud for processing.
This thesis designs and evaluates a \textbf{Fog Computing} system for monitoring
the power consumption of 40 smart homes, with an aggregate load of up to
\textbf{400 msg/s}, built on Apache Storm. The system is directly compared with a
\textbf{Monolithic} architecture (all processing on the Cloud) across six
experimental scenarios.

The measurements show that the Monolithic architecture saturates at around 20
homes (200 msg/s), where \code{Bolt\_split} reaches \textbf{169\%} capacity and
the complete latency rises to \textbf{5.28 s} at 400 msg/s. In contrast, the Fog
architecture with \textbf{8 gateways} keeps the cloud bolt capacity at
\textbf{0\%} at 400 msg/s, while the gateway processes tuples at sub-millisecond
latency (\textless 0.1 ms execute latency). A disk-backed Store-and-Forward
mechanism achieves \textbf{0\% data loss} during a 5-minute cloud outage and
fully recovers within \textbf{71--96 s}. WAN traffic is reduced by approximately
\textbf{86\%} through batched GZIP compression. The gateway is deployed on real
Raspberry Pi 3 hardware. The thesis contributes \textbf{eleven technical
improvements} over the Monolithic baseline, including an \textit{Energy
Optimizer} daemon that freezes idle gateways to bring their CPU usage to 0\%.

\vspace{0.8em}
\noindent\textbf{Keywords:} Fog Computing, Edge Computing, IoT, Apache Storm,
Smart Home, Store-and-Forward, Raspberry Pi, MQTT.

%% ---------- Mục lục, danh mục hình, bảng ----------
\tableofcontents
\listoffigures
\listoftables

%% ---------- Danh mục từ viết tắt ----------
\chapter*{DANH MỤC TỪ VIẾT TẮT}
\addcontentsline{toc}{chapter}{Danh mục từ viết tắt}
\thispagestyle{plain}
\renewcommand{\arraystretch}{1.3}
\begin{longtable}{p{0.22\textwidth} p{0.72\textwidth}}
\toprule
\textbf{Từ viết tắt} & \textbf{Diễn giải} \\
\midrule
\endhead
IoT      & Internet of Things --- Mạng lưới vạn vật kết nối \\
MQTT     & Message Queuing Telemetry Transport --- giao thức nhắn tin publish/subscribe \\
QoS      & Quality of Service --- mức bảo đảm gửi tin \\
WAN      & Wide Area Network --- mạng diện rộng \\
EC2      & Amazon Elastic Compute Cloud \\
JVM      & Java Virtual Machine --- máy ảo Java \\
JIT      & Just-In-Time compilation \\
JDBC     & Java Database Connectivity \\
GZIP     & GNU Zip --- thuật toán nén dữ liệu \\
EWMA     & Exponentially Weighted Moving Average --- trung bình trượt có trọng số mũ \\
RAM      & Random Access Memory \\
CPU      & Central Processing Unit \\
DB       & Database --- cơ sở dữ liệu \\
PK       & Primary Key --- khoá chính \\
NIST     & National Institute of Standards and Technology \\
KB       & Kịch bản (trong ngữ cảnh thực nghiệm) hoặc Kilobyte (đơn vị dữ liệu) \\
GW       & Gateway --- cổng biên Fog \\
EMA      & Exponential Moving Average \\
ZK       & ZooKeeper --- dịch vụ điều phối của Apache Storm \\
\bottomrule
\end{longtable}
\renewcommand{\arraystretch}{1.0}

%% ============================================================================
%%  PHẦN NỘI DUNG CHÍNH (đánh số Ả Rập)
%% ============================================================================
\pagenumbering{arabic}

%% ############################################################################
\chapter{MỞ ĐẦU}
\label{ch:modau}
%% ############################################################################

\section{Bối cảnh và động lực nghiên cứu}

Internet of Things (IoT) đang trải qua giai đoạn tăng trưởng bùng nổ. Theo nhiều
dự báo công nghiệp, số lượng thiết bị IoT kết nối toàn cầu sẽ vượt mốc
\textbf{50 tỷ thiết bị vào năm 2030}, trải rộng trên các lĩnh vực nhà thông minh,
y tế từ xa, nông nghiệp chính xác và thành phố thông minh. Mỗi thiết bị liên tục
phát sinh các luồng dữ liệu cảm biến với tần suất cao, tạo ra một khối lượng dữ
liệu khổng lồ cần được thu thập, xử lý và lưu trữ theo thời gian thực.

Mô hình điện toán truyền thống đẩy \textit{toàn bộ} dữ liệu thô từ thiết bị lên
Cloud để xử lý tập trung. Cách tiếp cận này tuy đơn giản nhưng bộc lộ ba hạn chế
nghiêm trọng khi quy mô tăng lên:

\begin{itemize}
  \item \textbf{Áp lực băng thông:} mỗi thông điệp thô được gửi riêng lẻ qua mạng
        diện rộng (WAN) khiến băng thông tiêu thụ tăng tuyến tính theo số lượng
        thiết bị. Với hàng trăm thiết bị, chi phí truyền dẫn và độ trễ mạng trở
        thành điểm nghẽn.
  \item \textbf{Độ trễ xử lý:} dữ liệu phải vượt qua một chặng đường dài
        (thiết bị $\to$ WAN $\to$ Cloud $\to$ xử lý $\to$ phản hồi), không phù hợp
        với các ứng dụng yêu cầu phản ứng tức thời như cảnh báo bất thường.
  \item \textbf{Phụ thuộc kết nối:} nếu đường truyền tới Cloud bị gián đoạn, toàn
        bộ pipeline ngừng hoạt động và dữ liệu trong thời gian mất kết nối bị mất
        hoàn toàn.
\end{itemize}

\textbf{Fog Computing} (điện toán sương mù) ra đời như một giải pháp cho các hạn
chế trên. Thay vì xử lý tập trung ở Cloud, Fog Computing bổ sung một tầng tính
toán trung gian đặt gần nguồn dữ liệu --- gọi là \textit{tầng sương mù} --- để
tiền xử lý, tổng hợp và lọc dữ liệu trước khi truyền phần tinh gọn lên Cloud.

Để nghiên cứu định lượng lợi ích của Fog Computing, đồ án này xây dựng một bài
toán cụ thể: \textbf{hệ thống giám sát tiêu thụ điện cho 40 ngôi nhà thông
minh}. Mỗi ngôi nhà phát các bản tin đo công suất với tốc độ \textbf{10 msg/s}
liên tục, tạo ra tổng tải tối đa \textbf{400 msg/s}. Hệ thống cần tổng hợp dữ
liệu theo nhiều cửa sổ thời gian và lưu trữ kết quả phục vụ trực quan hoá thời
gian thực.

\section{Vấn đề đặt ra}

Khi triển khai bài toán trên bằng kiến trúc \textbf{Monolithic} --- toàn bộ
pipeline xử lý (phân loại, tính trung bình, tổng hợp, dự báo) chạy trên một
cụm Apache Storm đặt ở Cloud --- hệ thống bộc lộ điểm bão hoà rõ rệt. Qua đo
lường thực nghiệm (\Cref{ch:ketqua-mono}), kiến trúc Monolithic đạt điểm bão hoà
tại khoảng \textbf{200 msg/s} (20 nhà): bolt phân loại \code{Bolt\_split} vượt
\textbf{169\%} capacity tại 400 msg/s và complete latency tăng vọt lên
\textbf{5,28 s}. Khi đó tỷ lệ ack giảm từ 3,4\% xuống còn 0,4\%, cho thấy hệ
thống không còn xử lý kịp tốc độ dữ liệu đến.

Câu hỏi đặt ra là: \textit{liệu một kiến trúc Fog Computing được thiết kế phù hợp
có thể vượt qua điểm nghẽn này, đồng thời cải thiện khả năng chịu lỗi và giảm tải
WAN hay không?} Và nếu có, cái giá phải trả về tài nguyên ở tầng biên là bao
nhiêu?

\section{Mục tiêu nghiên cứu}

Đồ án đặt ra bốn mục tiêu chính:

\begin{enumerate}
  \item Thiết kế và cài đặt một kiến trúc Fog Computing hoàn chỉnh trên nền tảng
        Apache Storm, với gateway biên chạy Storm LocalCluster và Cloud chạy cụm
        Storm phân tán.
  \item Đề xuất và hiện thực hoá các \textbf{cải tiến kỹ thuật} so với kiến trúc
        Monolithic, bao gồm tích luỹ gia tăng, nén GZIP theo lô, Store-and-Forward,
        phát hiện bất thường tại biên và tối ưu năng lượng.
  \item Đo lường và so sánh \textbf{tám chỉ số hiệu năng} trên \textbf{sáu kịch
        bản thực nghiệm} giữa hai kiến trúc.
  \item Triển khai và kiểm chứng gateway trên thiết bị nhúng tài nguyên thấp
        \textbf{Raspberry Pi 3}.
\end{enumerate}

\section{Phạm vi và giới hạn}

\textbf{Phạm vi:} đồ án tập trung vào quy mô 40 ngôi nhà, sử dụng dữ liệu mô phỏng
theo phong cách dataset REFIT, trên ngăn xếp công nghệ Apache Storm + MQTT +
MySQL + Prometheus + Grafana, triển khai trên Amazon EC2 (vùng ap-southeast-1) và
Raspberry Pi 3.

\textbf{Giới hạn:} đồ án không so sánh với các nền tảng xử lý luồng khác như Apache
Spark Streaming hay Apache Flink; không đo lường mức tiêu thụ điện năng phần cứng
của Raspberry Pi 3; và trong kịch bản 8 gateway, hệ thống dùng 7 container Docker
mô phỏng cộng với 1 Raspberry Pi 3 vật lý thực tế (chưa triển khai đủ 8 thiết bị
Pi vật lý).

\section{Đóng góp chính}

Đóng góp cốt lõi của đồ án là \textbf{mười cải tiến kỹ thuật} của kiến trúc Fog so
với Monolithic (bên cạnh một bản vá lỗi đo lường mang tính tiền đề), được trình
bày chi tiết trong \Cref{ch:thietke} và kiểm chứng bằng số liệu thực nghiệm trong
\Cref{ch:ketqua-fog,ch:sosanh}. Các cải tiến bao gồm: tích luỹ gia tăng thay cho
fan-out, nén GZIP theo lô, cơ chế Store-and-Forward chống mất dữ liệu, phát hiện
bất thường bằng EWMA z-score, định tuyến cảnh báo Slack theo từng nhà, bộ lập lịch
TagAwareScheduler, ghi idempotent bằng \code{REPLACE INTO}, đơn giản hoá topology
Cloud, và daemon Energy Optimizer tối ưu năng lượng tầng biên.

\section{Cấu trúc đồ án}

Phần còn lại của đồ án được tổ chức như sau. \textbf{\Cref{ch:cosoly}} trình bày cơ
sở lý thuyết về IoT, Fog Computing, Apache Storm và các công nghệ liên quan.
\textbf{\Cref{ch:thietke}} mô tả chi tiết thiết kế và cài đặt hai kiến trúc cùng 11
cải tiến kỹ thuật. \textbf{\Cref{ch:phuongphap}} trình bày phương pháp thực nghiệm
và sáu kịch bản đo. \textbf{\Cref{ch:ketqua-mono}} và \textbf{\Cref{ch:ketqua-fog}}
lần lượt trình bày kết quả của hệ thống Monolithic và Fog. \textbf{\Cref{ch:sosanh}}
so sánh và phân tích toàn diện hai kiến trúc. \textbf{\Cref{ch:ketluan}} kết luận
và đề xuất hướng phát triển.

%% ############################################################################
\chapter{CƠ SỞ LÝ THUYẾT}
\label{ch:cosoly}
%% ############################################################################

\section{Internet of Things và Smart Home}

\subsection{Kiến trúc IoT ba tầng}

Một hệ thống IoT điển hình được tổ chức thành ba tầng chức năng:

\begin{itemize}
  \item \textbf{Tầng cảm nhận (Perception):} gồm các cảm biến và bộ chấp hành thu
        thập dữ liệu vật lý (nhiệt độ, công suất điện, độ ẩm\dots) và chuyển đổi
        thành tín hiệu số.
  \item \textbf{Tầng mạng (Network):} chịu trách nhiệm truyền dữ liệu từ tầng cảm
        nhận lên tầng ứng dụng thông qua các giao thức như MQTT, CoAP, HTTP trên
        nền Wi-Fi, Ethernet hoặc mạng di động.
  \item \textbf{Tầng ứng dụng (Application):} xử lý, phân tích và trực quan hoá dữ
        liệu, cung cấp dịch vụ cho người dùng cuối.
\end{itemize}

\subsection{Giao thức MQTT}

MQTT (Message Queuing Telemetry Transport) là giao thức nhắn tin nhẹ theo mô hình
\textbf{publish/subscribe}, được thiết kế cho các thiết bị tài nguyên hạn chế và
mạng băng thông thấp, độ trễ cao. Các thành phần chính gồm \textit{broker} (máy
chủ trung gian), \textit{publisher} (bên phát) và \textit{subscriber} (bên nhận),
liên kết qua các \textit{topic}.

MQTT hỗ trợ ba mức Quality of Service (QoS): \textbf{QoS 0} (gửi nhiều nhất một
lần, không bảo đảm), \textbf{QoS 1} (gửi ít nhất một lần, có thể trùng lặp) và
\textbf{QoS 2} (gửi đúng một lần, bảo đảm cao nhất). Trong đồ án, các bản tin cảm
biến dùng QoS 0 để tối đa hoá thông lượng, trong khi tính idempotent ở Cloud đảm
bảo không sinh dữ liệu trùng.

\subsection{Dataset REFIT}

REFIT là tập dữ liệu công khai về điện năng tiêu thụ tại 20 hộ gia đình ở Anh,
trong đó công suất từng thiết bị được ghi nhận với chu kỳ khoảng 8 giây. Đồ án
mở rộng mô hình này lên 40 ngôi nhà bằng publisher mô phỏng, mỗi nhà phát bản tin
với tốc độ 10 msg/s nhằm tạo tải đủ lớn để khảo sát điểm bão hoà.

\section{Fog Computing}

\subsection{Định nghĩa}

Theo định nghĩa của NIST (SP 500-325), \textit{Fog Computing là một mô hình tính
toán phân tán nằm theo chiều ngang, cung cấp tài nguyên tính toán, lưu trữ, điều
khiển và mạng dọc theo liên tục từ Cloud đến thiết bị (Cloud-to-Thing)}. Fog
Computing không thay thế Cloud mà bổ sung một tầng trung gian giữa thiết bị biên
và trung tâm dữ liệu.

\subsection{So sánh Edge, Fog và Cloud}

\Cref{tab:edge-fog-cloud} so sánh ba mô hình theo các thuộc tính chính.

\begin{table}[H]
\centering
\caption{So sánh Edge, Fog và Cloud Computing}
\label{tab:edge-fog-cloud}
\begin{tabular}{p{0.24\textwidth} p{0.22\textwidth} p{0.22\textwidth} p{0.22\textwidth}}
\toprule
\textbf{Thuộc tính} & \textbf{Edge} & \textbf{Fog} & \textbf{Cloud} \\
\midrule
Vị trí          & Trên thiết bị      & Gần thiết bị (LAN)  & Trung tâm dữ liệu \\
Độ trễ          & Rất thấp (ms)      & Thấp (ms--vài chục) & Cao (vài trăm ms) \\
Năng lực tính toán & Hạn chế         & Trung bình          & Rất lớn \\
Lưu trữ         & Tạm thời          & Trung hạn           & Dài hạn, lớn \\
Băng thông WAN  & Không cần         & Giảm mạnh           & Tốn nhiều \\
Khả năng offline & Có               & Có                  & Không \\
\bottomrule
\end{tabular}
\end{table}

\subsection{Lợi ích và mô hình Store-and-Forward}

Fog Computing mang lại ba lợi ích cốt lõi: \textbf{giảm băng thông} (nhờ tổng hợp
và nén dữ liệu tại biên), \textbf{giảm độ trễ} (xử lý gần nguồn) và \textbf{khả
năng chịu lỗi offline} (tiếp tục hoạt động khi mất kết nối Cloud).

Khả năng offline được hiện thực hoá bằng mô hình \textbf{Store-and-Forward}: khi
đường truyền tới Cloud bị gián đoạn, tầng Fog \textit{lưu tạm} (store) các lô dữ
liệu đã tổng hợp xuống đĩa, và \textit{chuyển tiếp} (forward) chúng theo thứ tự
khi kết nối phục hồi. Nhờ đó, dữ liệu không bị mất trong thời gian mất kết nối ---
một đặc tính mà kiến trúc Monolithic không thể có.

\section{Apache Storm}

\subsection{Kiến trúc cụm}

Apache Storm là một hệ thống xử lý luồng dữ liệu phân tán, thời gian thực. Các
thành phần chính gồm:

\begin{itemize}
  \item \textbf{Nimbus:} nút điều phối trung tâm, phân phối mã và giao việc cho
        các Supervisor.
  \item \textbf{Supervisor:} chạy trên mỗi nút worker, khởi động và giám sát các
        tiến trình Worker theo lệnh của Nimbus.
  \item \textbf{Worker:} một tiến trình JVM thực thi một phần của topology.
  \item \textbf{Executor:} một luồng (thread) trong Worker, chạy một hoặc nhiều
        Task.
  \item \textbf{Task:} một thực thể của Spout hoặc Bolt thực thi logic xử lý.
\end{itemize}

\subsection{Các khái niệm cốt lõi}

\textbf{Topology} là đồ thị xử lý gồm các Spout và Bolt. \textbf{Spout} là nguồn
dữ liệu (đọc từ MQTT, Kafka\dots). \textbf{Bolt} là đơn vị xử lý (lọc, tổng hợp,
ghi DB\dots). \textbf{Stream} là luồng tuple vô hạn chảy giữa các thành phần.
\textbf{Grouping} quy định cách phân phối tuple từ một thành phần tới thành phần
kế tiếp (shuffle, fields, all\dots).

\subsection{Các metric quan trọng}

\Cref{tab:storm-metric} liệt kê các metric Storm dùng để đánh giá hiệu năng.

\begin{table}[H]
\centering
\caption{Các metric hiệu năng của Apache Storm}
\label{tab:storm-metric}
\begin{tabular}{p{0.26\textwidth} p{0.66\textwidth}}
\toprule
\textbf{Metric} & \textbf{Ý nghĩa} \\
\midrule
\code{bolt\_capacity}   & Tỷ lệ thời gian bolt bận xử lý; \textgreater 100\% nghĩa là bolt quá tải, không xử lý kịp tốc độ đến. \\
\code{complete\_latency}& Thời gian trung bình một tuple từ khi Spout phát ra đến khi được ack hoàn toàn. \\
\code{execute\_latency} & Thời gian thực thi trung bình một lần gọi \code{execute()}. \\
\code{acked / emitted}  & Số tuple được xác nhận / phát ra; tỷ lệ ack thấp cho thấy nghẽn cổ chai. \\
\bottomrule
\end{tabular}
\end{table}

\subsection{LocalCluster và Distributed cluster}

Apache Storm hỗ trợ hai chế độ chạy: \textbf{LocalCluster} mô phỏng toàn bộ cụm
trong một tiến trình JVM duy nhất (phù hợp cho thiết bị nhúng như Raspberry Pi),
và \textbf{Distributed cluster} chạy trên nhiều nút vật lý với Nimbus,
Supervisor và ZooKeeper riêng biệt (dùng cho Cloud). Trong đồ án, gateway Fog
dùng LocalCluster còn Cloud dùng cụm phân tán.

\subsection{Cơ chế đảm bảo xử lý và backpressure}
\label{sec:acking}

Storm bảo đảm mỗi tuple được xử lý ít nhất một lần (\textit{at-least-once}) thông
qua cơ chế \textbf{acking}: mỗi tuple do Spout phát ra được gán một ID, và Storm
theo dõi toàn bộ cây tuple sinh ra từ nó (anchoring). Khi tất cả các tuple con
được bolt gọi \code{ack()}, tuple gốc mới được coi là hoàn tất; nếu quá thời gian
\code{topology.message.timeout.secs} mà chưa hoàn tất, Spout sẽ phát lại tuple.
Chính cơ chế này định nghĩa \code{complete\_latency} --- khoảng thời gian từ khi
phát tuple đến khi cây tuple được ack đầy đủ.

Khi một bolt không xử lý kịp tốc độ tuple đến (capacity tiến tới và vượt 100\%),
hàng đợi nội bộ của nó đầy dần. Storm phản ứng bằng \textbf{backpressure}: tín
hiệu nghẽn lan ngược về phía Spout, buộc Spout giảm tốc độ phát (throttle). Đây
là lý do trong kịch bản Monolithic, khi \code{Bolt\_split} quá tải, tốc độ phát
của Spout bị kéo xuống còn khoảng 199--296 ops/s dù đầu vào là 400 msg/s. Hiểu rõ
acking và backpressure là điều kiện để diễn giải chính xác các số liệu capacity,
complete latency và spout emitted rate ở các chương kết quả.

\section{Các công nghệ liên quan}

\begin{itemize}
  \item \textbf{Docker và Docker Compose:} container hoá các dịch vụ thành các đơn
        vị độc lập, đảm bảo môi trường nhất quán và dễ triển khai. Docker Compose
        điều phối nhiều container qua một tệp khai báo.
  \item \textbf{Prometheus + Grafana:} Prometheus thu thập metric theo mô hình
        kéo (pull) với chu kỳ 5 giây; Grafana trực quan hoá metric thành biểu đồ
        thời gian thực.
  \item \textbf{MySQL 8.4:} hệ quản trị cơ sở dữ liệu quan hệ, lưu trữ dữ liệu
        tổng hợp dạng time-series.
  \item \textbf{GZIP:} thuật toán nén không mất mát, dùng để nén các lô dữ liệu
        trước khi truyền qua WAN.
  \item \textbf{Raspberry Pi 3:} máy tính nhúng một bo mạch với CPU ARM
        Cortex-A53 bốn nhân 1.2 GHz, 1 GB RAM, Wi-Fi tích hợp --- đại diện cho
        thiết bị biên tài nguyên thấp.
\end{itemize}

\section{Các chỉ số hiệu năng}

Đồ án đo lường tám chỉ số hiệu năng được định nghĩa đầy đủ trong
\Cref{tab:chiso}.

\begin{table}[H]
\centering
\caption{Tám chỉ số hiệu năng được đo lường}
\label{tab:chiso}
\small
\renewcommand{\arraystretch}{1.25}
\begin{tabular}{p{0.20\textwidth} c c p{0.20\textwidth} p{0.28\textwidth}}
\toprule
\textbf{Chỉ số} & \textbf{Ký hiệu} & \textbf{Đơn vị} & \textbf{Cách đo} & \textbf{Ý nghĩa} \\
\midrule
Storm Complete Latency & $L_{complete}$ & ms    & Prometheus (StormExporter) & Thời gian 1 tuple từ Spout $\to$ Ack \\
Bolt Capacity          & $C_{bolt}$     & \%    & Prometheus                 & \% thời gian bolt xử lý (\textgreater 100\% = quá tải) \\
Bolt Execute Latency   & $L_{exec}$     & ms    & Prometheus                 & Thời gian thực thi 1 lần \code{execute()} \\
End-to-End DB Latency  & $L_{e2e}$      & s     & MySQL query qua SSH        & Thời gian MQTT publish $\to$ ghi DB \\
WAN Traffic            & $B_{WAN}$      & KB/min& Prometheus (gateway)       & Băng thông gateway $\to$ cloud \\
Gateway RAM            & $M_{GW}$       & MB    & Docker stats               & Bộ nhớ mỗi gateway container \\
Gateway CPU            & $P_{GW}$       & \%    & Docker stats               & Tải CPU mỗi gateway \\
Store-Forward Queue    & $Q_{depth}$    & batch & Prometheus                 & Số batch đang chờ trong disk queue \\
\bottomrule
\end{tabular}
\renewcommand{\arraystretch}{1.0}
\end{table}

\section{Công trình liên quan}

Hướng nghiên cứu Fog/Edge Computing cho IoT đã thu hút nhiều công trình. Bonomi
và cộng sự \cite{bonomi2012} là những người đầu tiên hệ thống hoá vai trò của Fog
Computing trong IoT. NIST \cite{nist2018} và OpenFog Consortium
\cite{openfog2017} chuẩn hoá định nghĩa và kiến trúc tham chiếu. Shi và cộng sự
\cite{shi2016} phân tích Edge Computing như một hướng đi tất yếu khi dữ liệu biên
tăng nhanh. Yi và cộng sự \cite{yi2015} khảo sát các nền tảng và thách thức của
Fog. Trong lĩnh vực xử lý luồng, Apache Storm \cite{storm} và mô hình
Heron/Storm của Twitter \cite{toshniwal2014} được dùng rộng rãi cho phân tích
thời gian thực. Các nghiên cứu về MQTT \cite{mqtt} cho IoT và về cơ chế
store-and-forward \cite{psaras2018} cung cấp nền tảng cho thiết kế chịu lỗi của
đồ án. So với các công trình mô hình hoá topology gắn thẻ (tagged) trải dài giữa
Fog và Cloud trong \textit{cùng một} topology, đồ án này tách hẳn Fog và Cloud
thành \textit{hai} triển khai Storm độc lập, kết nối qua MQTT mang aggregate nén
--- một cách tiếp cận giảm triệt để lưu lượng WAN và cho phép biên hoạt động độc
lập khi Cloud mất kết nối.

%% ############################################################################
\chapter{THIẾT KẾ VÀ CÀI ĐẶT HỆ THỐNG}
\label{ch:thietke}
%% ############################################################################

\section{Kiến trúc tổng thể}

Đồ án xây dựng và so sánh hai kiến trúc xử lý cho cùng một bài toán giám sát điện
năng. \Cref{fig:arch-compare} đặt cạnh nhau hai kiến trúc Monolithic và Fog.

\begin{figure}[H]
\centering
\begin{tikzpicture}[
  node distance=0.6cm,
  box/.style={draw, rounded corners, align=center, minimum height=1cm, font=\small, fill=boxGray},
  cloud/.style={draw, rounded corners, align=center, minimum height=1cm, font=\small, fill=accentBlue!12},
  arr/.style={-{Stealth[length=2.5mm]}, thick}
]
%% --- Monolithic (trái) ---
\node[font=\bfseries] (m0) {(a) Monolithic};
\node[box, below=0.5cm of m0] (msrc) {Sensor /\\ Publisher};
\node[cloud, below=of msrc, text width=4cm] (mcloud)
  {\textbf{EC2 t3.large}\\ MQTT Broker + Storm\\ (Spout + Bolt\_split +\\ avg + sum + forecast)\\ + MySQL + Prometheus\\ + Grafana};
\draw[arr] (msrc) -- node[right, font=\scriptsize]{MQTT raw} (mcloud);

%% --- Fog (phải) ---
\node[font=\bfseries, right=4.5cm of m0] (f0) {(b) Fog};
\node[box, below=0.5cm of f0] (fsrc) {Sensor /\\ Publisher};
\node[box, below=of fsrc, text width=4cm] (fgw)
  {\textbf{Fog Gateway (Pi 3)}\\ Storm LocalCluster\\ Spout\_mqtt + Bolt\_ingest};
\node[box, right=1.8cm of fgw, text width=2.2cm, fill=fog1Orange!18] (qd)
  {Disk Queue\\ \code{queue.jsonl}};
\node[cloud, below=of fgw, text width=4cm] (fcloud)
  {\textbf{EC2 t3.small}\\ Storm (Spout +\\ Bolt\_cloudMerge)\\ + MySQL + Prometheus\\ + Grafana};
\draw[arr] (fsrc) -- node[right, font=\scriptsize]{MQTT raw} (fgw);
\draw[arr] (fgw) -- node[right, font=\scriptsize]{MQTT GZIP batch} (fcloud);
\draw[arr, dashed] (fgw) -- node[midway, above, font=\scriptsize]{S\&F} (qd);
\end{tikzpicture}
\caption{So sánh kiến trúc Monolithic (a) và Fog (b).}
\label{fig:arch-compare}
\end{figure}

Điểm khác biệt cốt lõi: trong kiến trúc Monolithic, mọi bản tin thô đều được đẩy
thẳng lên Cloud và toàn bộ pipeline tổng hợp chạy tập trung. Trong kiến trúc Fog,
gateway biên tổng hợp dữ liệu \textit{trước}, chỉ truyền các lô đã nén lên Cloud,
và có hàng đợi đĩa dự phòng cho tình huống mất kết nối.

\section{Kiến trúc Monolithic}

Topology Monolithic xử lý dữ liệu theo \textbf{tám cửa sổ thời gian}
(1, 5, 10, 15, 20, 30, 60, 120 phút). \Cref{fig:mono-topo} mô tả luồng xử lý.

\begin{figure}[H]
\centering
\begin{tikzpicture}[
  node distance=1.0cm and 1.4cm,
  sp/.style={draw, ellipse, align=center, font=\small, fill=accentBlue!12, minimum width=2.2cm},
  bo/.style={draw, rounded corners, align=center, font=\small, fill=boxGray, minimum height=0.9cm},
  arr/.style={-{Stealth[length=2.5mm]}, thick}
]
\node[sp] (spd) {Spout\_data};
\node[bo, right=of spd] (split) {Bolt\_split\\(phân loại theo\\ time-window)};
\node[bo, right=2.4cm of split, yshift=2.2cm] (avg) {Bolt\_avg\\(trung bình,\\ flush DB)};
\node[bo, right=2.4cm of split] (sum) {Bolt\_sum\\(tổng hợp\\ household)};
\node[bo, right=2.4cm of split, yshift=-2.2cm] (fc) {Bolt\_forecast\\(dự báo)};
\node[sp, below=2.6cm of spd] (spt) {Spout\_trigger};
\draw[arr] (spd) -- (split);
\draw[arr] (split) -- (avg);
\draw[arr] (split) -- (sum);
\draw[arr] (split) -- (fc);
\draw[arr] (spt) -| node[below, font=\scriptsize, pos=0.25]{kích hoạt flush định kỳ} (avg);
\end{tikzpicture}
\caption{Topology Monolithic: Spout\_data $\to$ Bolt\_split $\to$ \{avg, sum, forecast\}.}
\label{fig:mono-topo}
\end{figure}

Lược đồ cơ sở dữ liệu Monolithic là bảng \code{device\_data} với khoá chính phức
hợp \code{(house\_id, household\_id, device\_id, year, month, day, slice\_gap,
slice\_index)}, lưu giá trị trung bình theo từng cửa sổ thời gian.

\textbf{Điểm nghẽn:} \code{Bolt\_split} phải xử lý 100\% bản tin thô và phân loại
chúng vào tám cửa sổ thời gian cho từng thiết bị. Với cấu hình hiện tại, bolt này
không song song hoá hiệu quả, trở thành điểm nghẽn khi tải vượt 200 msg/s. Đây
chính là nguyên nhân bão hoà sẽ được phân tích định lượng trong
\Cref{ch:ketqua-mono}.

\section{Kiến trúc Fog --- Thiết kế từng thành phần}

\subsection{Fog Gateway (Raspberry Pi 3)}
\label{sec:fog-gw}

Tầng biên Fog được triển khai theo hai kịch bản:

\textbf{Kịch bản 8 gateway (KB1-A, KB2-A):} hệ thống gồm \textbf{7 container Docker}
chạy trên máy chủ phát triển và \textbf{1 Raspberry Pi 3 vật lý} chạy gateway
gw-08, phụ trách 5 ngôi nhà cuối cùng (house-35 đến house-39) --- tổng cộng 8
gateway, mỗi gateway phụ trách 5 nhà. Pi 3 chạy container \code{fog-gateway:latest},
kết nối Wi-Fi tới MQTT simulator và Cloud EC2.

\textbf{Kịch bản 1 gateway (KB1-B, KB2-B):} \textbf{một Raspberry Pi 3 duy nhất}
chạy toàn bộ gateway, phụ trách tất cả 40 nhà, thể hiện khả năng triển khai trên
thiết bị nhúng tài nguyên thấp.

\Cref{fig:deploy-gw} minh hoạ sơ đồ triển khai tầng biên.

\begin{figure}[H]
\centering
\begin{tikzpicture}[
  box/.style={draw, rounded corners, align=center, font=\small, fill=boxGray, minimum height=0.9cm},
  pi/.style={draw, rounded corners, align=center, font=\small, fill=fog1Orange!18, minimum height=0.9cm},
  cl/.style={draw, rounded corners, align=center, font=\small, fill=accentBlue!12, minimum height=0.9cm},
  arr/.style={-{Stealth[length=2.5mm]}, thick}
]
\node[box, text width=5cm] (docker)
  {\textbf{MacBook / Server}\\ 7 container Docker:\\ gw-01 \dots\ gw-07};
\node[pi, right=2.0cm of docker, text width=3.6cm] (pi)
  {\textbf{Raspberry Pi 3}\\ ARM Cortex-A53 1.2GHz\\ 1GB RAM, Docker\\ gw-08};
\node[cl, below=1.6cm of $(docker)!0.5!(pi)$, text width=4cm] (ec2)
  {\textbf{EC2 Cloud}\\ Storm + MySQL +\\ Prometheus + Grafana};
\draw[arr] (docker) -- node[left, font=\scriptsize, pos=0.6]{MQTT GZIP} (ec2);
\draw[arr] (pi) -- node[right, font=\scriptsize, pos=0.6]{MQTT GZIP} (ec2);
\end{tikzpicture}
\caption{Sơ đồ triển khai tầng biên kịch bản 8 gateway: 7 container Docker + 1 Raspberry Pi 3.}
\label{fig:deploy-gw}
\end{figure}

\subsection{Storm LocalCluster tại Gateway}

Topology gateway tối giản: \code{Spout\_mqtt} $\to$ \code{Bolt\_ingest}. Bolt
\code{Bolt\_ingest} thực hiện năm bước:

\begin{enumerate}
  \item Subscribe MQTT topic \code{iot-data}.
  \item Parse JSON payload \code{\{house\_id, household\_id, device\_id, value, event\_ts\}}.
  \item Cập nhật \textbf{5 accumulator} (\code{ConcurrentHashMap}) theo 5 cửa sổ
        thời gian 1, 5, 10, 15, 30 phút --- mỗi accumulator lưu cặp
        \textit{(count, sum)}, \textbf{không lưu dữ liệu thô}.
  \item Flush mỗi 60 giây: serialize toàn bộ accumulator $\to$ JSON $\to$ GZIP
        $\to$ MQTT publish lên topic \code{fog-data} của Cloud.
  \item \textbf{Store-and-Forward:} nếu publish thất bại, ghi lô vào tệp
        \code{queue.jsonl} trên đĩa.
\end{enumerate}

\Cref{fig:ingest-flow} mô tả luồng xử lý của \code{Bolt\_ingest}.

\begin{figure}[H]
\centering
\begin{tikzpicture}[
  node distance=0.85cm,
  io/.style={draw, trapezium, trapezium left angle=70, trapezium right angle=110, align=center, font=\scriptsize, fill=accentBlue!12, minimum width=2.4cm},
  proc/.style={draw, rounded corners, align=center, font=\scriptsize, fill=boxGray, minimum width=3.2cm, minimum height=0.8cm},
  dec/.style={draw, diamond, aspect=2, align=center, font=\scriptsize, fill=fog8Green!15, inner sep=1pt},
  qd/.style={draw, align=center, font=\scriptsize, fill=fog1Orange!18, minimum width=2.6cm, minimum height=0.8cm},
  arr/.style={-{Stealth[length=2.2mm]}, thick}
]
\node[io] (in) {Nhận MQTT tuple};
\node[proc, below=of in] (parse) {Parse JSON payload};
\node[proc, below=of parse] (acc) {Cập nhật 5 accumulator\\ (count++, sum += value)};
\node[dec, below=of acc] (timer) {Đến hạn\\ flush 60s?};
\node[proc, below=1.0cm of timer] (ser) {Serialize + GZIP batch};
\node[dec, below=of ser] (pub) {MQTT\\ publish OK?};
\node[proc, below=1.0cm of pub] (clear) {Xoá batch (thành công)};
\node[qd, right=1.6cm of pub] (queue) {append \code{queue.jsonl}};
\draw[arr] (in) -- (parse);
\draw[arr] (parse) -- (acc);
\draw[arr] (acc) -- (timer);
\draw[arr] (timer) -- node[right, font=\scriptsize]{Có} (ser);
\draw[arr] (timer.east) -- ++(2.6,0) |- node[pos=0.25, right, font=\scriptsize]{Chưa: chờ tuple} (in.east);
\draw[arr] (ser) -- (pub);
\draw[arr] (pub) -- node[right, font=\scriptsize]{Có} (clear);
\draw[arr] (pub) -- node[above, font=\scriptsize]{Không} (queue);
\end{tikzpicture}
\caption{Lưu đồ xử lý của \code{Bolt\_ingest}: parse $\to$ tích luỹ $\to$ flush $\to$ (thành công / queue).}
\label{fig:ingest-flow}
\end{figure}

\subsection{Cơ chế Store-and-Forward}
\label{sec:saf}

Cơ chế Store-and-Forward được mô hình hoá bằng máy trạng thái trong
\Cref{fig:saf-sm}. Khi MQTT publish thành công, hệ thống ở trạng thái
\textsc{Online} và xoá batch. Khi publish thất bại (Cloud offline), hệ thống
chuyển sang \textsc{Offline}, ghi batch vào \code{queue.jsonl}. Khi kết nối phục
hồi, hàm \code{drainQueue()} đọc tuần tự hàng đợi và publish lại từng entry, xoá
entry khi thành công.

\begin{figure}[H]
\centering
\begin{tikzpicture}[
  st/.style={draw, circle, align=center, font=\small, minimum size=1.8cm},
  arr/.style={-{Stealth[length=2.5mm]}, thick}
]
\node[st, fill=fog8Green!18] (on) {ONLINE};
\node[st, fill=monoRed!15, right=4.5cm of on] (off) {OFFLINE};
\draw[arr] (on) to[bend left=18] node[above, font=\scriptsize]{MQTT fail $\to$ append queue} (off);
\draw[arr] (off) to[bend left=18] node[below, font=\scriptsize]{kết nối lại $\to$ drainQueue()} (on);
\draw[arr] (on) to[loop above] node[above, font=\scriptsize]{MQTT ok $\to$ clear batch} (on);
\draw[arr] (off) to[loop above] node[above, font=\scriptsize]{tiếp tục append batch} (off);
\end{tikzpicture}
\caption{Máy trạng thái Store-and-Forward: ONLINE $\leftrightarrow$ OFFLINE.}
\label{fig:saf-sm}
\end{figure}

\textbf{Tính toán dung lượng hàng đợi:} mỗi gateway sinh ra
$5 \text{ windows} \times 5 \text{ flush cycles}$ trong 5 phút mất kết nối
$= 25$ batch. Với 8 gateway, tổng số batch tối đa là $8 \times 25 = 200$ batch ---
con số sẽ được kiểm chứng trong \Cref{ch:ketqua-fog}.

\subsection{Energy Optimizer --- Tự động tiết kiệm tài nguyên}
\label{sec:energy}

Một daemon Node.js (\code{tools/energy\_optimizer.js}) chạy song song với hệ
thống Fog, thực hiện \textbf{auto-scaling thông minh} tại tầng biên.

\textbf{Nguyên lý hoạt động:}
\begin{enumerate}
  \item Daemon subscribe topic MQTT \code{iot-data}, parse \code{houseId} từ mỗi
        bản tin.
  \item Cập nhật \code{gwLastSeen[gw] = Date.now()} cho gateway phụ trách nhà đó
        (ánh xạ: $\lfloor house\_id / 5 \rfloor + 1 = gw\_num$).
  \item Mỗi 3 giây, kiểm tra: nếu
        \code{Date.now() - gwLastSeen[gw] > 15000ms} $\to$ gateway nhàn rỗi $\to$
        \code{docker pause gw-0N}.
  \item Nếu traffic quay lại $\to$ \code{docker unpause gw-0N} để tiếp tục xử lý.
\end{enumerate}

\textbf{Tại sao dùng \code{docker pause} thay vì \code{docker stop}?} Lệnh
\code{docker stop} sẽ shutdown JVM, restart mất khoảng 10 giây và làm mất toàn bộ
accumulator state trong bộ nhớ. Ngược lại, \code{docker pause} dùng cơ chế
\textit{cgroups freezer} của Linux để đóng băng tiến trình: RAM được giữ nguyên,
CPU về 0\%, không mất state, và độ trễ unpause dưới 50 ms (xác nhận từ log:
\code{[ACTIVATE]} $\to$ \code{[SUCCESS]} trong khoảng 100 ms).

\textbf{Tiết kiệm tài nguyên:} một JVM gateway nhàn rỗi tiêu tốn khoảng 200 MB
RAM và 2--5\% CPU; khi bị pause, nó vẫn giữ 200 MB RAM nhưng CPU về \textbf{0\%}.
Trong kịch bản thực tế (ví dụ nấc h10 chỉ có 2 gateway active), 6 gateway nhàn rỗi
được pause, tiết kiệm khoảng 12--30\% CPU của máy chủ.

\Cref{fig:gw-life} mô tả vòng đời gateway. Lưu ý: kịch bản Fog-1GW không dùng
Energy Optimizer vì chỉ có một gateway --- không có gì để pause chọn lọc.

\begin{figure}[H]
\centering
\begin{tikzpicture}[
  st/.style={draw, rounded corners, align=center, font=\small, minimum width=2.6cm, minimum height=1.0cm},
  arr/.style={-{Stealth[length=2.5mm]}, thick}
]
\node[st, fill=fog8Green!18] (run) {RUNNING};
\node[st, fill=boxGray, right=4.5cm of run] (pause) {PAUSED\\ (cgroups freezer)};
\draw[arr] (run) to[bend left=20] node[above, font=\scriptsize]{nhàn rỗi \textgreater 15s $\to$ docker pause} (pause);
\draw[arr] (pause) to[bend left=20] node[below, font=\scriptsize]{traffic quay lại $\to$ unpause \textless 50ms} (run);
\end{tikzpicture}
\caption{Máy trạng thái vòng đời gateway dưới điều khiển của Energy Optimizer.}
\label{fig:gw-life}
\end{figure}

\subsection{Cloud Storm Topology (Fog)}

Topology Cloud trong kiến trúc Fog đơn giản hơn Monolithic rất nhiều --- chỉ gồm
một bolt: \code{Spout\_mqtt} (subscribe \code{fog-data}) $\to$
\code{Bolt\_cloudMerge} (REPLACE INTO MySQL). \Cref{fig:fog-topo} thể hiện toàn bộ
pipeline Fog từ biên đến Cloud.

\begin{figure}[H]
\centering
\begin{tikzpicture}[
  node distance=1.0cm,
  sp/.style={draw, ellipse, align=center, font=\scriptsize, fill=accentBlue!12, minimum width=1.8cm},
  bo/.style={draw, rounded corners, align=center, font=\scriptsize, fill=boxGray, minimum height=0.8cm},
  db/.style={draw, cylinder, shape border rotate=90, aspect=0.3, align=center, font=\scriptsize, fill=fog8Green!12, minimum height=1cm},
  arr/.style={-{Stealth[length=2.2mm]}, thick}
]
\node[sp] (sm) {Spout\_mqtt};
\node[bo, right=of sm] (bi) {Bolt\_ingest};
\node[bo, right=1.4cm of bi, fill=accentBlue!10] (sm2) {Spout\_mqtt\\ (cloud)};
\node[bo, right=of sm2] (cm) {Bolt\_cloudMerge};
\node[db, right=of cm] (my) {MySQL};
\node[font=\scriptsize, above=0.1cm of sm] {\textbf{Gateway (Pi 3)}};
\node[font=\scriptsize, above=0.1cm of cm] {\textbf{Cloud (EC2)}};
\draw[arr] (sm) -- (bi);
\draw[arr] (bi) -- node[above, font=\scriptsize, align=center]{MQTT\\GZIP} (sm2);
\draw[arr] (sm2) -- (cm);
\draw[arr] (cm) -- node[above, font=\scriptsize, align=center]{REPLACE\\INTO} (my);
\end{tikzpicture}
\caption{Pipeline Fog đầy đủ: từ gateway (Spout\_mqtt $\to$ Bolt\_ingest) đến Cloud (Spout\_mqtt $\to$ Bolt\_cloudMerge $\to$ MySQL).}
\label{fig:fog-topo}
\end{figure}

\code{Bolt\_cloudMerge.execute()} thực hiện: (1) nhận lô MQTT đã nén GZIP,
(2) giải nén và parse các bản ghi JSON đã tổng hợp, (3) thực thi một lệnh
\code{REPLACE INTO fog\_device\_data (...) VALUES (...)}, (4) đảm bảo tính
idempotent --- cùng khoá sẽ được ghi đè, không tạo dòng trùng. Lược đồ Cloud là
bảng \code{fog\_device\_data} với các cột \code{(house\_id, household\_id,
device\_id, slice\_gap, slice\_index, value, count, updatedAt)}.

\section{Khắc phục lỗi đo lường metric Prometheus (bug fix)}
\label{sec:prometheus-fix}

\textit{Phần này mô tả một bản vá lỗi (bug fix), \textbf{không phải} một cải tiến
kiến trúc.} Nó được trình bày riêng ở đây vì là điều kiện tiên quyết để toàn bộ số
liệu đo lường trong \Cref{ch:ketqua-fog} đáng tin cậy; các cải tiến kỹ thuật thực
sự của kiến trúc Fog được trình bày trong \Cref{sec:improvements}.

Phiên đo đầu tiên (June 14--16) gặp một lỗi phần mềm làm sai lệch toàn bộ số liệu
gateway. Khi Storm LocalCluster restart một bolt task trong cùng JVM (lần gọi
\code{prepare()} thứ hai), lệnh \code{new PrometheusMetricsServer(port)} ném
\code{IOException: Address already in use}, khiến biến \code{metrics} bị gán
\code{null} và mọi counter (\code{incrementTuplesProcessed()}\dots) bị bỏ qua
\textit{một cách âm thầm} (silent failure) --- không có exception lan ra ngoài.
Hệ quả là các gateway gw-03 đến gw-08 không báo bất kỳ metric nào, khiến
Prometheus/Grafana chỉ ``nhìn thấy'' 2/8 gateway hoạt động, dẫn tới kết luận sai
rằng tải dồn hết vào hai gateway đầu.

Nguyên nhân gốc là vòng đời cổng (port) không được quản lý xuyên suốt các lần
\code{prepare()}. Giải pháp dùng một \code{ConcurrentHashMap<Integer,
PrometheusMetricsServer>} ở phạm vi \code{static} làm sổ đăng ký (registry) toàn
JVM, kết hợp \code{computeIfAbsent}: nếu cổng đã được bind thì tái sử dụng instance
cũ thay vì tạo mới. Sau khi vá, cả 8 gateway báo metric chính xác và phiên đo được
thực hiện lại ngày 2026-06-28. Bài học rút ra là: trong môi trường đo lường, một
\textit{silent failure} ở tầng quan trắc nguy hiểm hơn một lỗi gây dừng chương
trình, vì nó tạo ra số liệu trông hợp lệ nhưng sai.

\begin{lstlisting}[language=Java,caption={Singleton METRIC\_REGISTRY với computeIfAbsent}]
private static final ConcurrentHashMap<Integer, PrometheusMetricsServer>
    METRIC_REGISTRY = new ConcurrentHashMap<>();

// Trong prepare(): reuse server neu cong da bind, tranh IOException
this.metrics = METRIC_REGISTRY.computeIfAbsent(port,
    p -> new PrometheusMetricsServer(p));
\end{lstlisting}

\section{Mười cải tiến kỹ thuật so với Monolithic}
\label{sec:improvements}

Đây là phần đóng góp kỹ thuật trọng tâm của đồ án. Kiến trúc Fog không chỉ ``dời''
việc xử lý xuống biên mà tái thiết kế lại từng thành phần của pipeline. Phần này
phân tích sâu \textbf{mười cải tiến} ở cấp độ kiến trúc: với mỗi cải tiến, đồ án
trình bày \textit{vấn đề gốc} của Monolithic, \textit{cơ chế} giải quyết của Fog
(kèm sơ đồ luồng hoặc mã giả khi cần), \textit{sự đánh đổi} đi kèm, và
\textit{tác động đo được} trên số liệu thực nghiệm. \Cref{tab:improvements} tóm
tắt trước toàn bộ; các tiểu mục sau phân tích chi tiết.

{\scriptsize
\renewcommand{\arraystretch}{1.25}
\begin{longtable}{c p{0.16\textwidth} p{0.24\textwidth} p{0.24\textwidth} p{0.18\textwidth}}
\caption{Mười cải tiến kỹ thuật của kiến trúc Fog so với Monolithic}
\label{tab:improvements}\\
\toprule
\textbf{\#} & \textbf{Thành phần} & \textbf{Vấn đề Monolithic} & \textbf{Giải pháp Fog} & \textbf{Tác động} \\
\midrule
\endfirsthead
\toprule
\textbf{\#} & \textbf{Thành phần} & \textbf{Vấn đề Monolithic} & \textbf{Giải pháp Fog} & \textbf{Tác động} \\
\midrule
\endhead
1  & Xử lý raw data & Bolt\_split nhận 100\% raw (400 msg/s) $\to$ 169\% capacity & Tích luỹ gia tăng trong memory, không fan-out & GW bolt capacity \textbf{\textless 1\%} \\
2  & 5 time windows & 8 windows qua 4 bolt $\to$ $O(n{\times}8)$ tuple & 5 windows trong 1 bolt, 1 accumulator dict & Không có traffic nội bộ \\
3  & GZIP batch & Raw JSON từng message lên cloud & 60s buffer $\to$ GZIP batch & WAN giảm \textbf{$\sim$86\%} \\
4  & Store-and-Forward & Cloud offline = 100\% mất dữ liệu & \code{queue.jsonl} + drainQueue() & \textbf{0\%} loss, 96s recovery \\
5  & EWMA Z-score & Không có phát hiện bất thường & EWMA $\mu/\sigma$ mỗi device, alert nếu $|z|>3$ & Phát hiện anomaly tại biên \\
6  & Slack routing & Không có cảnh báo & Webhook riêng mỗi nhà & Alert \textless 1s \\
7  & TagAwareScheduler & Storm không bảo đảm affinity & Gán executor đúng supervisor theo tag & Ổn định assignment \\
8  & REPLACE INTO & INSERT $\to$ rủi ro trùng khi retry & \code{REPLACE INTO} + composite PK & \textbf{0} dòng trùng \\
9  & Cloud topology & 4 bolt + 8 window streams & 1 bolt (Bolt\_cloudMerge) & Cloud capacity \textbf{0\%} tại h40 \\
10 & Energy Optimizer & Idle JVM tiêu 2--5\% CPU & Daemon \code{docker pause} GW idle \textgreater 15s, unpause \textless 50ms & CPU idle \textbf{0\%} \\
\bottomrule
\end{longtable}
\renewcommand{\arraystretch}{1.0}
}

\subsection{Cải tiến 1 --- Tích luỹ gia tăng thay cho fan-out}

\textbf{Vấn đề gốc.} Trong kiến trúc Monolithic, mỗi bản tin thô đến
\code{Bolt\_split} bị \textit{nhân bản} (fan-out) thành nhiều tuple --- một tuple
cho mỗi cửa sổ thời gian và mỗi nhánh tổng hợp. Với tám cửa sổ, một bản tin đầu
vào có thể sinh ra hàng chục tuple nội bộ. Mỗi tuple đó lại phải được serialize,
đẩy vào hàng đợi giữa các executor, được ack riêng và đôi khi tái phân phối qua
mạng. Hệ quả là khối lượng tuple lưu thông trong topology tăng theo cấp số nhân so
với số bản tin đầu vào: quan sát thực tế cho tỷ lệ \code{emitted/acked} vượt
\textbf{50:1}. Đây chính là nguyên nhân gốc rễ khiến \code{Bolt\_split} đạt
\textbf{169\%} capacity tại 400 msg/s.

\textbf{Cơ chế Fog.} Kiến trúc Fog thay mô hình fan-out bằng \textbf{tích luỹ gia
tăng} (incremental accumulation). Thay vì phát tuple, \code{Bolt\_ingest} chỉ cập
nhật \textit{tại chỗ} cặp \textit{(count, sum)} trong bộ nhớ cho mỗi cửa sổ
(\code{count++}, \code{sum += value}). \Cref{fig:imp-fanout} đối chiếu trực tiếp
hai mô hình: bên Monolithic, một bản tin lan thành ``cây tuple''; bên Fog, một bản
tin chỉ chạm vào vài ô nhớ rồi kết thúc.

\begin{figure}[H]
\centering
\begin{tikzpicture}[
  box/.style={draw, rounded corners, align=center, font=\scriptsize, minimum height=0.8cm},
  arr/.style={-{Stealth[length=2mm]}, thick}
]
%% (a) Monolithic fan-out
\node[font=\bfseries\scriptsize] (ta) {(a) Monolithic --- fan-out};
\node[box, fill=accentBlue!12, below=0.35cm of ta] (am) {1 bản\\ tin};
\node[box, fill=monoRed!22, right=0.9cm of am] (as) {Bolt\_split};
\node[box, fill=monoRed!10, right=0.9cm of as, text width=1.7cm] (at) {$\times 8$ tuple\\ (8 cửa sổ)};
\node[box, fill=boxGray, right=0.9cm of at, text width=1.7cm] (ab) {4 tầng bolt\\ (avg/sum/\dots)};
\draw[arr] (am) -- (as);
\draw[arr] (as) -- node[above, font=\tiny]{nhân bản} (at);
\draw[arr] (at) -- (ab);
%% (b) Fog accumulate
\node[font=\bfseries\scriptsize, below=1.7cm of ta, anchor=west, xshift=-0.0cm] (tb) {(b) Fog --- tích luỹ gia tăng};
\node[box, fill=accentBlue!12, below=0.35cm of tb] (bm) {1 bản\\ tin};
\node[box, fill=fog8Green!20, right=0.9cm of bm, text width=2.2cm] (bi) {Bolt\_ingest:\\ count++, sum+=value};
\node[box, fill=fog8Green!8, right=0.9cm of bi, text width=1.7cm] (bz) {0 tuple\\ mạng};
\draw[arr] (bm) -- (bi);
\draw[arr] (bi) -- node[above, font=\tiny]{cập nhật ô nhớ} (bz);
\end{tikzpicture}
\caption{So sánh fan-out (Monolithic) và tích luỹ gia tăng (Fog) cho một bản tin đầu vào.}
\label{fig:imp-fanout}
\end{figure}

\textbf{Đánh đổi và tác động.} Toàn bộ chi phí xử lý một bản tin ở Fog chỉ là vài
phép gán trên \code{ConcurrentHashMap}, nên execute latency của gateway luôn dưới
\textbf{0,1 ms} và bolt capacity giữ dưới \textbf{1\%} ngay cả ở tải đỉnh
(\Cref{tab:fog-kb1a}). Cái giá phải trả là gateway chỉ lưu \textit{trạng thái tổng
hợp} chứ không giữ chuỗi sự kiện gốc --- nghĩa là không thể truy vấn lại từng số
đọc thô sau khi đã tổng hợp. Với bài toán giám sát điện năng theo cửa sổ thời
gian, đây là đánh đổi hoàn toàn chấp nhận được, thậm chí có lợi vì giảm cả lưu trữ
lẫn băng thông.

\subsection{Cải tiến 2 --- Năm cửa sổ trong một bolt}

\textbf{Vấn đề gốc.} Monolithic phân tán tám cửa sổ thời gian qua bốn tầng bolt
(\code{split} $\to$ \code{avg} $\to$ \code{sum} $\to$ \code{forecast}), tạo ra độ
phức tạp $O(n \times 8)$ tuple với $n$ là số bản tin. Việc trải logic qua nhiều
tầng còn làm phát sinh \textit{phụ thuộc thứ tự} và chi phí đồng bộ giữa các bolt.

\textbf{Cơ chế Fog.} Kiến trúc Fog gộp toàn bộ năm cửa sổ (1, 5, 10, 15, 30 phút)
vào \textbf{một từ điển accumulator duy nhất} trong \code{Bolt\_ingest}. Khoá là
bộ \textit{(device, window)}, giá trị là cặp \textit{(count, sum)}.
\Cref{fig:imp-acc} minh hoạ: một số đọc 254\,W của thiết bị $d$ chỉ cập nhật năm ô
tương ứng với năm cửa sổ, mỗi ô là một phép cộng dồn $O(1)$.

\begin{figure}[H]
\centering
\begin{tikzpicture}[
  cell/.style={draw, align=center, font=\scriptsize, minimum width=3.0cm, minimum height=0.65cm, fill=fog8Green!10},
  arr/.style={-{Stealth[length=2mm]}, thick}
]
\node[draw, rounded corners, fill=accentBlue!12, align=center, font=\scriptsize] (r) {Số đọc mới\\ $d=$Plug0, $x=254$\,W};
\node[cell, right=1.4cm of r, yshift=1.3cm] (w1) {(d,\,1m) $\to$ (count, sum)};
\node[cell, below=0.12cm of w1] (w2) {(d,\,5m) $\to$ (count, sum)};
\node[cell, below=0.12cm of w2] (w3) {(d,\,10m) $\to$ (count, sum)};
\node[cell, below=0.12cm of w3] (w4) {(d,\,15m) $\to$ (count, sum)};
\node[cell, below=0.12cm of w4] (w5) {(d,\,30m) $\to$ (count, sum)};
\node[draw, dashed, fit=(w1)(w5), inner sep=4pt, label={[font=\tiny]above:ConcurrentHashMap}] (box) {};
\foreach \w in {w1,w2,w3,w4,w5} \draw[arr] (r.east) -- (\w.west);
\end{tikzpicture}
\caption{Một số đọc cập nhật đồng thời năm ô accumulator (mỗi cửa sổ một ô), không sinh tuple.}
\label{fig:imp-acc}
\end{figure}

\textbf{Tác động.} Độ phức tạp giảm từ $O(n \times 8)$ tuple mạng (Monolithic)
xuống $O(1)$ thao tác bộ nhớ trên mỗi bản tin (Fog). Vì không còn tuple nội bộ
giữa các bolt, gateway không phát sinh \textit{bất kỳ} lưu lượng mạng nội bộ nào
--- một điều kiện then chốt giúp thiết bị Raspberry Pi 3 tài nguyên thấp vẫn xử lý
được 50 msg/s với headroom lớn.

\subsection{Cải tiến 3 --- Nén GZIP theo lô}

\textbf{Vấn đề gốc.} Monolithic đẩy từng bản tin JSON thô lên Cloud ngay khi nhận,
nên lưu lượng WAN tỷ lệ thuận với tốc độ bản tin (ước tính 2.400 KB/min tại
400 msg/s) và bùng nổ mỗi khi có burst. Băng thông WAN vừa đắt vừa khó dự đoán.

\textbf{Cơ chế Fog.} Kiến trúc Fog kết hợp hai kỹ thuật, minh hoạ ở
\Cref{fig:imp-gzip}: \textbf{(1) traffic shaping} --- mỗi gateway gom dữ liệu trong
cửa sổ 60 giây rồi mới flush một lần, nên Cloud không bao giờ thấy burst; và
\textbf{(2) nén GZIP} --- dữ liệu tổng hợp dạng JSON có tính lặp cao (cùng tên
trường, cùng cấu trúc lặp lại) nên đạt tỷ lệ nén rất tốt.

\begin{figure}[H]
\centering
\begin{tikzpicture}[
  box/.style={draw, rounded corners, align=center, font=\scriptsize, minimum height=0.9cm, text width=2.2cm},
  arr/.style={-{Stealth[length=2.2mm]}, thick}
]
\node[box, fill=accentBlue!12] (a) {Buffer 60s\\ ($\sim$600 số đọc)};
\node[box, fill=boxGray, right=1.0cm of a] (b) {Serialize JSON\\ aggregate};
\node[box, fill=fog1Orange!18, right=1.0cm of b] (c) {GZIP\\ $\to \sim$8,3 KB/batch};
\node[box, fill=fog8Green!18, right=1.0cm of c] (d) {MQTT publish\\ topic fog-data};
\draw[arr] (a) -- (b);
\draw[arr] (b) -- (c);
\draw[arr] (c) -- (d);
\end{tikzpicture}
\caption{Đường ống flush của gateway: gom 60s $\to$ serialize $\to$ nén GZIP $\to$ publish.}
\label{fig:imp-gzip}
\end{figure}

\textbf{Tác động.} Lưu lượng WAN tại h40 chỉ còn \textbf{332,8 KB/min} --- giảm
khoảng \textbf{86\%} so với raw streaming. Quan trọng không kém con số tuyệt đối là
\textit{hình dạng} của lưu lượng: WAN gần như phẳng theo thời gian thay vì nhấp nhô
theo burst, giúp việc cấp phát băng thông WAN trở nên dễ dự đoán. Một quan sát tinh
tế (\Cref{ch:sosanh}) là cấu hình 1 gateway lại tốn WAN \textit{cao hơn} 8 gateway
19\%, vì batch lớn chứa nhiều khoá khác nhau làm giảm tỷ lệ nén --- cho thấy hiệu
quả nén phụ thuộc cấu trúc batch chứ không đơn thuần số lượng nguồn.

\subsection{Cải tiến 4 --- Store-and-Forward chống mất dữ liệu}

\textbf{Vấn đề gốc.} Trong Monolithic, nếu đường truyền tới Cloud (hoặc chính
Cloud) gián đoạn, không có nơi nào lưu giữ dữ liệu --- toàn bộ số đọc trong thời
gian gián đoạn mất vĩnh viễn (100\% data loss).

\textbf{Cơ chế Fog.} Mỗi chu kỳ flush, gateway thử publish lô lên Cloud; nếu thất
bại, lô được ghi nối tiếp xuống tệp \code{queue.jsonl} trên đĩa
(\Cref{lst:flush}). Khi kết nối phục hồi, hàm \code{drainQueue()} đọc tuần tự hàng
đợi và publish lại từng lô, chỉ xoá entry khi gửi thành công. Máy trạng thái tổng
quát đã trình bày ở \Cref{fig:saf-sm}; \Cref{fig:imp-drain} chi tiết hoá logic
\textit{drain} sau khi Cloud online trở lại.

\begin{lstlisting}[language=Java,caption={Bolt\_ingest.flush() với Store-and-Forward},label={lst:flush}]
void flush() {
  byte[] batch = gzip(serialize(accumulators)); // count+sum snapshot
  try {
    mqtt.publish("fog-data", batch);            // duong di binh thuong
  } catch (MqttException e) {
    appendToDiskQueue("queue.jsonl", batch);    // store-and-forward
  }
  resetAccumulators();
}
\end{lstlisting}

\begin{figure}[H]
\centering
\begin{tikzpicture}[
  node distance=0.7cm,
  io/.style={draw, rounded corners, align=center, font=\scriptsize, fill=accentBlue!12, minimum height=0.7cm},
  dec/.style={draw, diamond, aspect=2.2, align=center, font=\scriptsize, fill=fog8Green!15, inner sep=1pt},
  proc/.style={draw, align=center, font=\scriptsize, fill=boxGray, minimum height=0.7cm},
  arr/.style={-{Stealth[length=2.2mm]}, thick}
]
\node[io] (start) {Cloud online};
\node[dec, below=of start] (empty) {queue rỗng?};
\node[proc, below=1.0cm of empty] (read) {Đọc lô kế tiếp\\ từ queue.jsonl};
\node[dec, below=of read] (pub) {publish OK?};
\node[proc, below=1.0cm of pub] (del) {Xoá entry};
\node[io, right=1.8cm of empty] (done) {Hoàn tất drain};
\node[proc, right=1.8cm of pub] (stop) {Dừng,\\ chờ lần sau};
\draw[arr] (start) -- (empty);
\draw[arr] (empty) -- node[left, font=\tiny]{Không} (read);
\draw[arr] (empty) -- node[above, font=\tiny]{Có} (done);
\draw[arr] (read) -- (pub);
\draw[arr] (pub) -- node[left, font=\tiny]{OK} (del);
\draw[arr] (pub) -- node[above, font=\tiny]{lỗi} (stop);
\draw[arr] (del) -| (empty);
\end{tikzpicture}
\caption{Lưu đồ \code{drainQueue()}: xả tuần tự hàng đợi đĩa sau khi Cloud phục hồi.}
\label{fig:imp-drain}
\end{figure}

\textbf{Tác động.} Cơ chế này đạt \textbf{0\% mất dữ liệu} trong các thử nghiệm
outage 5 phút (\Cref{ch:ketqua-fog}), với thời gian phục hồi 71--96 s. Nhờ ghi
xuống đĩa (không chỉ giữ trong RAM), dữ liệu còn sống sót qua cả sự cố restart
gateway.

\subsection{Cải tiến 5 --- Phát hiện bất thường EWMA Z-score tại biên}

\textbf{Vấn đề gốc.} Trong kiến trúc cloud-centric, muốn phát hiện bất thường phải
chờ dữ liệu lên Cloud rồi mới phân tích --- vừa trễ, vừa tốn băng thông để truyền
cả dữ liệu bình thường lẫn bất thường.

\textbf{Cơ chế Fog.} \code{Bolt\_ingest} duy trì trung bình trượt có trọng số mũ
($\mu$) và phương sai trượt ($\sigma^2$) cho \textit{mỗi thiết bị}, cập nhật theo
hệ số làm mượt $\alpha = 0{,}1$:
\begin{align}
\mu_t &= (1-\alpha)\,\mu_{t-1} + \alpha\, x_t \\
\sigma^2_t &= (1-\alpha)\,\sigma^2_{t-1} + \alpha\,(x_t - \mu_{t-1})^2 \\
z_t &= \frac{x_t - \mu_t}{\sigma_t}
\end{align}
Nếu $|z_t| > 3$ (quy tắc ba-sigma), một cảnh báo được phát sinh ngay tại gateway.
\Cref{fig:imp-ewma} mô tả luồng xử lý. EWMA được chọn thay cho trung bình cửa sổ
trượt thông thường vì chỉ cần lưu \textit{hai số thực} cho mỗi thiết bị thay vì
toàn bộ chuỗi giá trị --- phù hợp ràng buộc bộ nhớ của Raspberry Pi 3 --- đồng thời
tự thích nghi với sự thay đổi chậm của nền tải điện.

\begin{figure}[H]
\centering
\begin{tikzpicture}[
  node distance=0.7cm,
  io/.style={draw, rounded corners, align=center, font=\scriptsize, fill=accentBlue!12, minimum height=0.7cm},
  proc/.style={draw, align=center, font=\scriptsize, fill=boxGray, minimum height=0.7cm},
  dec/.style={draw, diamond, aspect=2.2, align=center, font=\scriptsize, fill=fog8Green!15, inner sep=1pt},
  al/.style={draw, align=center, font=\scriptsize, fill=monoRed!15, minimum height=0.7cm},
  arr/.style={-{Stealth[length=2.2mm]}, thick}
]
\node[io] (x) {Số đọc $x_t$};
\node[proc, right=of x] (z) {Tính $z=\dfrac{x_t-\mu}{\sigma}$};
\node[dec, right=of z] (chk) {$|z|>3$?};
\node[al, right=of chk] (alert) {Tạo STORM-ALERT\\ $\to$ route Slack};
\node[proc, below=1.0cm of chk] (upd) {Cập nhật $\mu,\sigma$ (EWMA)};
\draw[arr] (x) -- (z);
\draw[arr] (z) -- (chk);
\draw[arr] (chk) -- node[above, font=\tiny]{Có} (alert);
\draw[arr] (chk) -- node[right, font=\tiny]{Không} (upd);
\draw[arr] (alert) |- (upd);
\end{tikzpicture}
\caption{Lưu đồ phát hiện bất thường EWMA z-score trong \code{Bolt\_ingest}.}
\label{fig:imp-ewma}
\end{figure}

\textbf{Tác động.} Việc phát hiện diễn ra ngay tại biên với chi phí $O(1)$ bộ nhớ
mỗi thiết bị, không phụ thuộc Cloud. Cảnh báo thực tế sinh ra từ cơ chế này được
minh hoạ ở \Cref{fig:slack-alerts} (mục tiếp theo): mỗi cảnh báo ghi rõ giá trị
đọc, baseline EWMA và độ lệch theo sigma, giúp người vận hành hiểu \textit{vì sao}
một số đọc bị coi là bất thường.

\subsection{Cải tiến 6 --- Định tuyến cảnh báo Slack theo từng nhà}

\textbf{Vấn đề gốc.} Phát hiện được bất thường là chưa đủ; cảnh báo phải đến đúng
người phụ trách đúng ngôi nhà, nếu không sẽ bị nhiễu loạn giữa hàng chục nhà.

\textbf{Cơ chế Fog.} Mỗi ngôi nhà được cấu hình một webhook Slack riêng (lưu trong
\code{.env.gateway}). Khi có cảnh báo, \code{Bolt\_ingest} tra cứu webhook theo
\code{house\_id} và POST thẳng tới kênh tương ứng (\Cref{fig:imp-slack-route}).
Toàn bộ chuỗi phát hiện $\to$ định tuyến $\to$ gửi diễn ra tại gateway nên độ trễ
cảnh báo dưới một giây và \textit{không} phụ thuộc kết nối Cloud --- ngay cả khi
Cloud offline, cảnh báo an toàn vẫn được gửi (miễn gateway còn truy cập Internet).

\begin{figure}[H]
\centering
\begin{tikzpicture}[
  box/.style={draw, rounded corners, align=center, font=\scriptsize, minimum height=0.7cm},
  arr/.style={-{Stealth[length=2.2mm]}, thick}
]
\node[box, fill=monoRed!15] (a) {Anomaly\\ house\_id=1};
\node[box, fill=boxGray, right=1.0cm of a, text width=2.4cm] (m) {Bảng webhook\\ house\_id $\to$ URL};
\node[box, fill=fog8Green!18, right=1.0cm of m] (c1) {\#house-1};
\node[box, fill=fog8Green!8, above=0.25cm of c1] (c0) {\#house-0};
\node[box, fill=fog8Green!8, below=0.25cm of c1] (c3) {\#house-3};
\draw[arr] (a) -- (m);
\draw[arr] (m) -- node[above, font=\tiny]{lookup} (c1);
\end{tikzpicture}
\caption{Định tuyến cảnh báo: tra webhook theo \code{house\_id} rồi gửi đúng kênh Slack.}
\label{fig:imp-slack-route}
\end{figure}

\Cref{fig:slack-alerts} là ảnh chụp thực tế kênh \code{\#house-1} nhận các cảnh báo
\code{STORM-ALERT} do gateway \code{gw-single} phát hiện bằng bộ \code{Z\_SCORE}.

\begin{figure}[H]
\centering
\IfFileExists{images/fog/slack_alerts.png}
  {\includegraphics[width=0.92\textwidth]{images/fog/slack_alerts.png}}
  {\fbox{\parbox[c][5cm][c]{0.9\textwidth}{\centering\itshape [Đặt ảnh chụp Slack tại \texttt{images/fog/slack\_alerts.png}]}}}
\caption{Cảnh báo bất thường thực tế trên Slack: mỗi cảnh báo ghi rõ giá trị đọc, baseline EWMA, độ lệch sigma và lý do (gateway gw-single, detector Z\_SCORE).}
\label{fig:slack-alerts}
\end{figure}

\subsection{Cải tiến 7 --- TagAwareScheduler (gán executor theo nhãn)}

\textbf{Vấn đề gốc.} Bộ lập lịch mặc định của Storm không bảo đảm một executor cụ
thể được đặt trên một supervisor cụ thể; bolt tổng hợp và bolt ghi DB có thể bị
phân tán ngẫu nhiên qua nhiều nút, làm tăng số chặng mạng và khiến mỗi lần triển
khai cho ra một cách phân bổ khác nhau (khó tái lập thực nghiệm).

\textbf{Cơ chế Fog.} Đồ án hiện thực hoá một \textbf{TagAwareScheduler} tuỳ biến:
mỗi thành phần topology được gán nhãn (\code{tier=gateway} hoặc \code{tier=cloud}),
và bộ lập lịch chỉ ánh xạ executor đến supervisor có nhãn tương ứng
(\Cref{fig:imp-tag}). Nhờ đó \code{Bolt\_cloudMerge} luôn được ghim trên nút Cloud
kề MySQL, còn \code{Bolt\_ingest} ở lại gateway.

\begin{figure}[H]
\centering
\begin{tikzpicture}[
  box/.style={draw, rounded corners, align=center, font=\scriptsize, minimum height=0.7cm},
  arr/.style={-{Stealth[length=2.2mm]}, thick}
]
\node[box, fill=fog8Green!18] (e1) {Executor\\ tier=gateway};
\node[box, fill=accentBlue!14, below=0.5cm of e1] (e2) {Executor\\ tier=cloud};
\node[box, fill=boxGray, right=1.4cm of e1, yshift=-0.55cm, text width=2.0cm] (sch) {TagAware\\ Scheduler};
\node[box, fill=fog8Green!10, right=1.6cm of sch, yshift=0.55cm] (s1) {Supervisor\\ (gateway)};
\node[box, fill=accentBlue!8, right=1.6cm of sch, yshift=-0.55cm] (s2) {Supervisor\\ (cloud)};
\draw[arr] (e1) -- (sch);
\draw[arr] (e2) -- (sch);
\draw[arr] (sch) -- node[above, font=\tiny]{match tag} (s1);
\draw[arr] (sch) -- node[below, font=\tiny]{match tag} (s2);
\end{tikzpicture}
\caption{TagAwareScheduler ánh xạ executor đến supervisor cùng nhãn (tier).}
\label{fig:imp-tag}
\end{figure}

\textbf{Tác động.} Việc gán cố định giảm số chặng mạng giữa các tầng và làm cho
phân bổ topology \textit{ổn định, lặp lại được} giữa các lần triển khai --- điều
kiện cần để các phép đo ở \Cref{ch:ketqua-fog} có thể so sánh công bằng.

\subsection{Cải tiến 8 --- Ghi idempotent với REPLACE INTO}

\textbf{Vấn đề gốc.} Trong cơ chế Store-and-Forward, một lô có thể được gửi lại
nhiều lần khi kết nối chập chờn (ví dụ publish bị timeout nhưng thực tế đã tới
Cloud). Nếu Cloud dùng \code{INSERT}, mỗi lần gửi lại tạo thêm một dòng trùng và
làm sai lệch số liệu tổng hợp.

\textbf{Cơ chế Fog.} Cloud dùng \code{REPLACE INTO} với khoá chính phức hợp
\code{(house\_id, household\_id, device\_id, slice\_gap, slice\_index)}: cùng một
khoá sẽ \textit{ghi đè} thay vì chèn mới. Kết hợp với ngữ nghĩa \textit{cumulative}
(gateway luôn gửi tổng tích luỹ chứ không phải số gia), thao tác ghi trở nên
\textbf{idempotent} hoàn toàn. \Cref{fig:imp-idem} minh hoạ: dù lô ``total=100''
được gửi hai lần, hàng trong DB vẫn là 100 chứ không cộng thành 200.

\begin{figure}[H]
\centering
\begin{tikzpicture}[
  box/.style={draw, rounded corners, align=center, font=\scriptsize, minimum height=0.7cm},
  db/.style={draw, cylinder, shape border rotate=90, aspect=0.3, align=center, font=\scriptsize, fill=fog8Green!12, minimum height=0.9cm},
  arr/.style={-{Stealth[length=2.2mm]}, thick}
]
\node[box, fill=accentBlue!12] (g) {Gateway};
\node[db, right=3.2cm of g] (d) {DB row\\ =100};
\draw[arr] (g.north east) ++(0,0.1) -- node[above, font=\tiny]{gửi lần 1: total=100} ([yshift=0.5cm]d.west);
\draw[arr] (g.south east) ++(0,-0.1) -- node[below, font=\tiny]{gửi lại: total=100 (REPLACE)} ([yshift=-0.5cm]d.west);
\end{tikzpicture}
\caption{Ghi idempotent: gửi lại cùng lô cumulative không làm sai số liệu (REPLACE giữ giá trị 100).}
\label{fig:imp-idem}
\end{figure}

\textbf{Tác động.} Thực nghiệm xác nhận \textbf{0 dòng trùng} sau quá trình drain
recovery, bảo đảm tính đúng đắn của dữ liệu ngay cả khi mạng không ổn định.

\subsection{Cải tiến 9 --- Đơn giản hoá topology Cloud}

\textbf{Vấn đề gốc.} Ở Monolithic, Cloud phải gánh toàn bộ pipeline: phân loại,
nhân bản theo cửa sổ, tính trung bình, tổng hợp household và dự báo --- bốn tầng
bolt với tám luồng cửa sổ, chính là nơi sinh ra nút thắt \code{Bolt\_split}.

\textbf{Cơ chế Fog.} Vì gateway đã làm hết phần tổng hợp, topology Cloud được rút
gọn còn \textbf{một bolt duy nhất} (\code{Bolt\_cloudMerge}) chỉ làm nhiệm vụ giải
nén lô và \code{REPLACE INTO} DB. \Cref{fig:imp-cloudtopo} đối chiếu hai topology.

\begin{figure}[H]
\centering
\begin{tikzpicture}[
  box/.style={draw, rounded corners, align=center, font=\tiny, minimum height=0.6cm},
  arr/.style={-{Stealth[length=1.8mm]}, semithick}
]
%% Mono
\node[font=\bfseries\scriptsize] (mt) {(a) Cloud Monolithic};
\node[box, fill=accentBlue!12, below=0.3cm of mt] (ms) {Spout};
\node[box, fill=monoRed!22, right=0.5cm of ms] (msp) {split};
\node[box, fill=boxGray, right=0.5cm of msp] (mavg) {avg};
\node[box, fill=boxGray, right=0.4cm of mavg] (msum) {sum};
\node[box, fill=boxGray, right=0.4cm of msum] (mfc) {forecast};
\draw[arr] (ms)--(msp); \draw[arr] (msp)--(mavg); \draw[arr] (mavg)--(msum); \draw[arr] (msum)--(mfc);
%% Fog
\node[font=\bfseries\scriptsize, below=1.3cm of mt, anchor=west] (ft) {(b) Cloud Fog};
\node[box, fill=accentBlue!12, below=0.3cm of ft] (fs) {Spout};
\node[box, fill=fog8Green!20, right=0.5cm of fs, text width=1.8cm] (fm) {Bolt\_cloudMerge};
\node[box, fill=fog8Green!10, right=0.6cm of fm] (fdb) {MySQL};
\draw[arr] (fs)--(fm); \draw[arr] (fm)--node[above,font=\tiny]{REPLACE}(fdb);
\end{tikzpicture}
\caption{Topology Cloud: Monolithic bốn tầng bolt (a) so với Fog một bolt (b).}
\label{fig:imp-cloudtopo}
\end{figure}

\textbf{Tác động.} Đây là lý do trực tiếp khiến cloud bolt capacity về \textbf{0\%}
ở trạng thái ổn định tại h40 trong kịch bản 8 gateway: lượng công việc trên mỗi
đơn vị thời gian ở Cloud giảm hàng chục lần. Mặt trái cần lưu ý là khi \textit{một}
gateway gửi quá nhiều bản ghi trong một lô (cấu hình 1 gateway, 200 dòng/batch),
chính \code{Bolt\_cloudMerge} lại trở thành nút thắt mới (318,9\% tại h40) --- vấn
đề được phân tích ở \Cref{ch:sosanh}.

\subsection{Cải tiến 10 --- Energy Optimizer daemon}

\textbf{Vấn đề gốc.} Trong cấu hình nhiều gateway, ở các giai đoạn tải thấp nhiều
gateway không có nhà nào active nhưng JVM của chúng vẫn tiêu tốn 2--5\% CPU mỗi cái
một cách lãng phí.

\textbf{Cơ chế Fog.} Một daemon Node.js (\code{energy\_optimizer.js}) giám sát
traffic MQTT theo thời gian thực và tự động ``đóng băng'' gateway nhàn rỗi
(\Cref{fig:imp-energy-loop}): cứ mỗi 3 giây, nếu một gateway không nhận bản tin nào
quá 15 giây thì daemon gọi \code{docker pause}; khi traffic quay lại thì
\code{docker unpause}. Máy trạng thái vòng đời gateway đã trình bày ở
\Cref{fig:gw-life}. Điểm tinh tế là dùng \code{pause}/\code{unpause} (cgroups
freezer) thay cho \code{stop}/\code{start}: nó đưa CPU về 0\% nhưng \textit{giữ
nguyên} accumulator state trong RAM và unpause trong dưới 50 ms, tránh được độ trễ
khởi động lại JVM ($\sim$10 s) lẫn mất state.

\begin{figure}[H]
\centering
\begin{tikzpicture}[
  node distance=0.65cm,
  io/.style={draw, rounded corners, align=center, font=\scriptsize, fill=accentBlue!12, minimum height=0.7cm},
  proc/.style={draw, align=center, font=\scriptsize, fill=boxGray, minimum height=0.7cm},
  dec/.style={draw, diamond, aspect=2.4, align=center, font=\scriptsize, fill=fog8Green!15, inner sep=1pt},
  act/.style={draw, align=center, font=\scriptsize, fill=fog1Orange!18, minimum height=0.7cm},
  arr/.style={-{Stealth[length=2.2mm]}, thick}
]
\node[io] (sub) {Subscribe iot-data};
\node[proc, below=of sub] (upd) {Cập nhật gwLastSeen[gw]};
\node[dec, below=of upd] (idle) {idle \textgreater 15s?};
\node[act, right=1.8cm of idle] (pause) {docker pause gw};
\node[dec, below=1.0cm of idle] (traf) {traffic\\ quay lại?};
\node[act, right=1.8cm of traf] (unp) {docker unpause gw};
\draw[arr] (sub) -- (upd);
\draw[arr] (upd) -- (idle);
\draw[arr] (idle) -- node[above, font=\tiny]{Có} (pause);
\draw[arr] (idle) -- node[right, font=\tiny]{Không} (traf);
\draw[arr] (traf) -- node[above, font=\tiny]{Có (\textless 50ms)} (unp);
\draw[arr] (traf.west) -- ++(-1.3,0) |- node[pos=0.25, left, font=\tiny]{lặp mỗi 3s} (sub.west);
\end{tikzpicture}
\caption{Vòng lặp giám sát của Energy Optimizer: pause gateway nhàn rỗi, unpause khi traffic quay lại.}
\label{fig:imp-energy-loop}
\end{figure}

\textbf{Tác động.} CPU của gateway nhàn rỗi về \textbf{0\%} (từ 2--5\%). Cơ chế đặc
biệt hiệu quả ở giai đoạn tải thấp --- ví dụ nấc h10 chỉ có 2/8 gateway active,
sáu gateway còn lại được pause --- tiết kiệm 12--30\% CPU máy chủ
(\Cref{fig:energy}). Vì chỉ có ý nghĩa khi tồn tại nhiều gateway, cải tiến này chỉ
áp dụng cho cấu hình 8 gateway.

\section{Sơ đồ triển khai phân tầng}

\Cref{fig:deploy-tiers} trình bày sơ đồ triển khai năm tầng của hệ thống Fog.

\begin{figure}[H]
\centering
\begin{tikzpicture}[
  tier/.style={draw, rounded corners, align=center, font=\small, minimum width=11cm, minimum height=1.0cm},
  arr/.style={-{Stealth[length=2.5mm]}, very thick}
]
\node[tier, fill=boxGray] (t1) {\textbf{Tầng 1 --- Sensors/Publishers:} 40 houses $\times$ simulator (MQTT raw 10 msg/s)};
\node[tier, fill=fog1Orange!16, below=0.7cm of t1] (t2) {\textbf{Tầng 2 --- Fog Edge:} Raspberry Pi 3 (+ 7 container Docker nếu chế độ 8-GW)};
\node[tier, fill=accentBlue!10, below=0.7cm of t2] (t3) {\textbf{Tầng 3 --- WAN:} MQTT GZIP batch (60s/flush)};
\node[tier, fill=fog8Green!14, below=0.7cm of t3] (t4) {\textbf{Tầng 4 --- AWS EC2 (ap-southeast-1):} Nimbus + Supervisor + ZK + MySQL + Prometheus + Grafana};
\node[tier, fill=boxGray, below=0.7cm of t4] (t5) {\textbf{Tầng 5 --- Web Dashboard} (port 9000)};
\draw[arr] (t1) -- (t2);
\draw[arr] (t2) -- (t3);
\draw[arr] (t3) -- (t4);
\draw[arr] (t4) -- (t5);
\end{tikzpicture}
\caption{Sơ đồ triển khai phân tầng của hệ thống Fog Computing.}
\label{fig:deploy-tiers}
\end{figure}

\section{Cấu hình phần cứng và phần mềm}

\Cref{tab:config} so sánh môi trường triển khai của ba hệ thống.

\begin{table}[H]
\centering
\caption{Cấu hình phần cứng và phần mềm ba hệ thống}
\label{tab:config}
\small
\begin{tabular}{p{0.22\textwidth} p{0.24\textwidth} p{0.24\textwidth} p{0.16\textwidth}}
\toprule
\textbf{Thành phần} & \textbf{Monolithic} & \textbf{Fog (8 GW)} & \textbf{Fog (1 GW)} \\
\midrule
Cloud instance & EC2 t3.large (2vCPU, 8GB) & EC2 t3.small (2vCPU, 2GB) & EC2 t3.small \\
Edge hardware  & ---                       & 7$\times$ Docker + 1$\times$ Pi 3 & 1$\times$ Pi 3 \\
Apache Storm   & 2.1.0 (distributed)       & LocalCluster             & LocalCluster \\
MySQL          & 8.4.2                     & 8.4.2                    & 8.4.2 \\
MQTT Broker    & mosquitto                 & mosquitto                & mosquitto \\
Prometheus     & 2.54.1                    & 2.54.1                   & 2.54.1 \\
Grafana        & 11.1.5                    & 11.1.5                   & 11.1.5 \\
\bottomrule
\end{tabular}
\end{table}

%% ############################################################################
\chapter{PHƯƠNG PHÁP THỰC NGHIỆM}
\label{ch:phuongphap}
%% ############################################################################

\section{Tổng quan sáu kịch bản đo}

Đồ án thực hiện sáu kịch bản đo, tóm tắt trong \Cref{tab:scenarios}.

\begin{table}[H]
\centering
\caption{Tổng quan sáu kịch bản thực nghiệm}
\label{tab:scenarios}
\small
\begin{tabular}{p{0.18\textwidth} c p{0.18\textwidth} p{0.20\textwidth} c}
\toprule
\textbf{Kịch bản} & \textbf{Ký hiệu} & \textbf{Hệ thống} & \textbf{Mục tiêu} & \textbf{Ngày đo} \\
\midrule
Ramp Monolithic   & MONO-KB1  & Monolithic            & Điểm bão hoà        & 2026-06-18 \\
Steady Monolithic & MONO-KB2  & Monolithic            & Headline metrics    & 2026-06-18 \\
Ramp Fog 8 GW     & FOG-KB1-A & Fog 8 GW (7 Docker+Pi)& Scalability phân tán& 2026-06-28 \\
Ramp Fog 1 GW     & FOG-KB1-B & Fog 1 GW (1 Pi 3)     & Scalability đơn     & 2026-06-28 \\
Recovery Fog 8 GW & FOG-KB2-A & Fog 8 GW              & Store-and-Forward   & 2026-06-28 \\
Recovery Fog 1 GW & FOG-KB2-B & Fog 1 GW              & Store-and-Forward   & 2026-06-28 \\
\bottomrule
\end{tabular}
\end{table}

\section{Dataset và Publisher}

Dữ liệu mô phỏng theo phong cách dataset REFIT (40 ngôi nhà, mỗi thiết bị đo công
suất với chu kỳ khoảng 8 giây). Publisher viết bằng Node.js, phát bản tin MQTT
với trường \code{event\_ts = Date.now()}, tốc độ \code{SPEED=10} msg/s mỗi nhà,
trên topic \code{iot-data}, QoS = 0.

\textbf{Kịch bản ramp} (scalability) tăng tải theo 5 nấc, mỗi nấc 360 giây:
\code{h01 $\to$ h05 $\to$ h10 $\to$ h20 $\to$ h40} (tương ứng 10, 50, 100, 200,
400 msg/s). \textbf{Kịch bản steady} giữ tải cố định 40 nhà trong 1.800 giây.

\section{Metrics thu thập và công cụ}

\Cref{tab:metric-tool} liệt kê các metric, nguồn, script và tần suất thu thập.

\begin{table}[H]
\centering
\caption{Metrics, nguồn và công cụ thu thập}
\label{tab:metric-tool}
\small
\begin{tabular}{p{0.26\textwidth} p{0.24\textwidth} p{0.22\textwidth} p{0.14\textwidth}}
\toprule
\textbf{Metric} & \textbf{Nguồn} & \textbf{Script} & \textbf{Tần suất} \\
\midrule
Storm Complete Latency & Prometheus StormExporter & \code{observe\_ramp*.sh} & Cuối mỗi nấc \\
Bolt Capacity          & Prometheus               & \code{observe\_ramp*.sh} & Cuối mỗi nấc \\
Gateway Exec Latency   & Prometheus (fog\_gateway\_*) & \code{observe\_ramp*.sh} & Cuối mỗi nấc \\
WAN KB/min             & Prometheus (flush\_total, publish\_size) & \code{observe\_ramp*.sh} & Cuối mỗi nấc \\
CPU/RAM                & Docker stats $\to$ python & \code{observe\_ramp*.sh} & Mỗi 12s \\
End-to-end Latency     & MySQL SSH query           & \code{latency\_report.sh} & Sau mỗi kịch bản \\
Queue Depth            & Prometheus (store\_queue\_size) & \code{kb2\_offline\_recovery.sh} & Mỗi 5s \\
\bottomrule
\end{tabular}
\end{table}

\subsection{Truy vấn Prometheus và SQL}

Các metric Storm và gateway được lấy trực tiếp từ Prometheus bằng PromQL. Ví dụ,
cloud bolt capacity được truy vấn theo nhãn topology và component, còn gateway
execute latency được tính bằng tỷ số giữa tổng thời gian và số lần gọi trong cửa
sổ một phút:

\begin{lstlisting}[caption={Hai truy vấn PromQL tiêu biểu}]
# Cloud bolt capacity (gia tri thap phan, 1.0 = 100%)
storm_bolt_capacity{topology="fog-cloud", component="Bolt_cloudMerge"}

# Gateway execute latency trung binh trong 1 phut (ms)
rate(fog_gateway_execute_latency_sum[1m])
  / rate(fog_gateway_execute_latency_count[1m])
\end{lstlisting}

End-to-end latency được đo bằng truy vấn SQL trực tiếp trên MySQL qua SSH, lấy
hiệu giữa thời điểm ghi DB (\code{updatedAt}) và thời điểm sự kiện gốc
(\code{event\_ts}), sau đó tính các phân vị p50/p95/p99:

\begin{lstlisting}[language=SQL,caption={Truy vấn SQL đo end-to-end latency}]
SELECT
  TIMESTAMPDIFF(SECOND, FROM_UNIXTIME(event_ts/1000), updatedAt) AS lat_s
FROM fog_device_data
WHERE updatedAt >= NOW() - INTERVAL 30 MINUTE
ORDER BY lat_s;
-- p50/p95/p99 duoc tinh tu phan phoi lat_s o tang script
\end{lstlisting}

Như đã lưu ý, do \code{updatedAt} luôn bằng \code{NOW()} của lần ghi cuối nên giá
trị \code{lat\_s} là artifact đo lường (xem \Cref{ch:sosanh} và
Phụ lục~\ref{app:artifact}).

\section{Script tự động hoá}

Ba script chính điều khiển toàn bộ quá trình đo:
\begin{itemize}
  \item \code{tools/observe\_ramp.sh} --- chạy ramp scalability cho Fog, tự động
        tăng tải và chụp snapshot metric cuối mỗi nấc.
  \item \code{tools/kb2\_offline\_recovery.sh} --- test offline recovery: warmup
        15 phút $\to$ tắt cloud-mqtt 5 phút $\to$ bật lại $\to$ đo recovery time.
  \item \code{tools/observe\_ramp\_mono.sh} --- chạy ramp scalability cho
        Monolithic.
\end{itemize}

\section{Hạn chế và biện pháp kiểm soát}

Để bảo đảm tính tin cậy, mỗi kịch bản được chạy lại nếu xuất hiện outlier, và
Docker được reset giữa các kịch bản để loại bỏ trạng thái tồn dư.

\textbf{Hạn chế đo lường End-to-End Latency:} do \code{REPLACE INTO} dùng
\code{NOW()} nên cột \code{updatedAt} luôn bằng thời điểm ghi \textit{cuối cùng},
không phải lần ghi đầu tiên. Vì vậy các giá trị End-to-End latency cao (401--655 s)
là \textbf{artifact thiết kế}, không phải độ trễ thực. Độ trễ ghi lần đầu thực tế
xấp xỉ 60--65 giây (chu kỳ flush + thời gian xử lý Cloud). Vấn đề này được phân
tích sâu trong \Cref{ch:sosanh} và Phụ lục~\ref{app:artifact}.

%% ############################################################################
\chapter{KẾT QUẢ HỆ THỐNG MONOLITHIC}
\label{ch:ketqua-mono}
%% ############################################################################

\section{MONO-KB1 --- Ramp Scalability}

\Cref{tab:mono-kb1} trình bày số liệu ramp của hệ thống Monolithic.

\begin{table}[H]
\centering
\caption{Số liệu MONO-KB1 (ramp scalability)}
\label{tab:mono-kb1}
\small
\begin{tabular}{c c c c c c c}
\toprule
\textbf{Số nhà} & \textbf{Msg/s} & \textbf{Complete Lat.} & \textbf{Bolt Split Cap.} & \textbf{Acked/600s} & \textbf{Emitted/600s} & \textbf{Ack Ratio} \\
\midrule
1  & 10  & 55 ms              & 45,5\%            & 2.840  & 82.460    & 3,4\% \\
5  & 50  & 97 ms              & 11,2\%            & 13.973 & 438.034   & 3,2\% \\
10 & 100 & 290 ms             & 3,9\%             & 30.326 & 1.057.009 & 2,9\% \\
20 & 200 & \textbf{1.111 ms}  & \textbf{89,8\%}   & 22.422 & 1.493.779 & 1,5\% \\
40 & 400 & \textbf{5.282 ms}  & \textbf{169\%}    & 6.873  & 1.547.121 & 0,4\% \\
\bottomrule
\end{tabular}
\end{table}

\Cref{fig:mono-cl} cho thấy complete latency tăng theo cấp số nhân khi số nhà
tăng, vượt ngưỡng degradation 1.000 ms ngay tại 20 nhà.

\begin{figure}[H]
\centering
\begin{tikzpicture}
\begin{axis}[
  width=0.82\textwidth, height=6.5cm,
  ybar, bar width=16pt,
  ymode=log, log basis y=10,
  ymin=10, ymax=10000,
  symbolic x coords={1,5,10,20,40},
  xtick=data,
  xlabel={Số nhà}, ylabel={Complete Latency (ms, log)},
  nodes near coords, point meta=explicit symbolic,
  nodes near coords align={vertical}, every node near coord/.append style={font=\scriptsize},
  enlarge x limits=0.15,
]
\addplot[fill=accentBlue!70] coordinates {(1,55)[55] (5,97)[97] (10,290)[290] (20,1111)[1.111] (40,5282)[5.282]};
\draw[dashed, monoRed, very thick] (axis cs:1,1000) -- (axis cs:40,1000)
  node[pos=0.5, above, font=\scriptsize, monoRed]{Ngưỡng degradation 1000 ms};
\end{axis}
\end{tikzpicture}
\caption{MONO-KB1: Storm Complete Latency theo số nhà (thang log).}
\label{fig:mono-cl}
\end{figure}

\Cref{fig:mono-cap} kết hợp Bolt Split Capacity (cột) và Ack Ratio (đường), làm
rõ điểm bão hoà tại 20 nhà.

\begin{figure}[H]
\centering
\begin{tikzpicture}
\begin{axis}[
  width=0.82\textwidth, height=6.5cm,
  axis y line*=left,
  ybar, bar width=16pt,
  symbolic x coords={1,5,10,20,40}, xtick=data,
  xlabel={Số nhà}, ylabel={Bolt Split Capacity (\%)},
  ymin=0, ymax=180, enlarge x limits=0.15,
  legend style={at={(0.02,0.98)}, anchor=north west, font=\scriptsize},
]
\addplot[fill=monoRed!70] coordinates {(1,45.5) (5,11.2) (10,3.9) (20,89.8) (40,169)};
\addlegendentry{Bolt Split Capacity}
\draw[dashed, black] (axis cs:1,100) -- (axis cs:40,100);
\end{axis}
\begin{axis}[
  width=0.82\textwidth, height=6.5cm,
  axis y line*=right, axis x line=none,
  symbolic x coords={1,5,10,20,40},
  ylabel={Ack Ratio (\%)}, ymin=0, ymax=4,
  legend style={at={(0.02,0.86)}, anchor=north west, font=\scriptsize},
]
\addplot[fog1Orange, very thick, mark=*] coordinates {(1,3.4) (5,3.2) (10,2.9) (20,1.5) (40,0.4)};
\addlegendentry{Ack Ratio}
\end{axis}
\end{tikzpicture}
\caption{MONO-KB1: Bolt Split Capacity (cột) và Ack Ratio (đường). Điểm bão hoà tại 20 nhà.}
\label{fig:mono-cap}
\end{figure}

\Cref{fig:mono-grafana-kb1} là ảnh chụp Grafana xác nhận các xu hướng trên.

\begin{figure}[H]
\centering
\begin{subfigure}{0.49\textwidth}
  \centering\grafana{images/mono/s1_ramp_bolt_capacity.png}
  \caption{Bolt capacity}
\end{subfigure}\hfill
\begin{subfigure}{0.49\textwidth}
  \centering\grafana{images/mono/s1_ramp_complete_latency.png}
  \caption{Complete latency}
\end{subfigure}
\caption{MONO-KB1 (Grafana): bolt capacity và complete latency theo ramp.}
\label{fig:mono-grafana-kb1}
\end{figure}

\section{Phân tích bottleneck Monolithic}

Tại nấc h40, \code{Bolt\_split} nhận toàn bộ 400 msg/s và phải phân loại vào tám
cửa sổ thời gian cho nhiều thiết bị. Với capacity 169\%, throughput thực chỉ đạt
khoảng $400 \times 1/1{,}69 \approx 237$ msg/s, dẫn tới tích luỹ
$400 - 237 = 163$ msg/s trong hàng đợi. Sau 360 giây của nấc h40, ước tính có
khoảng $163 \times 360 \approx 58.680$ bản tin chờ xử lý trong queue, gây timeout
và drop. \Cref{fig:mono-bottleneck} minh hoạ dòng nghẽn này.

\begin{figure}[H]
\centering
\begin{tikzpicture}[
  node distance=1.2cm,
  bo/.style={draw, rounded corners, align=center, font=\scriptsize, minimum height=0.9cm},
  arr/.style={-{Stealth[length=2.5mm]}, thick}
]
\node[bo, fill=accentBlue!12] (sp) {Spout\\ 400/s};
\node[bo, fill=monoRed!25, right=of sp] (split) {Bolt\_split\\ 169\% cap.};
\node[bo, fill=fog1Orange!20, right=of split] (q) {Queue\\ overflow};
\node[bo, fill=boxGray, right=of q] (avg) {Bolt\_avg\\ 8,9\%};
\node[bo, fill=monoRed!15, right=of avg] (drop) {timeout /\\ drop};
\draw[arr] (sp) -- (split);
\draw[arr] (split) -- node[above, font=\scriptsize]{237/s} (q);
\draw[arr] (q) -- (avg);
\draw[arr] (avg) -- (drop);
\end{tikzpicture}
\caption{Dòng nghẽn Monolithic tại h40: Bolt\_split 169\% capacity gây tràn hàng đợi.}
\label{fig:mono-bottleneck}
\end{figure}

\section{MONO-KB2 --- Steady State 40 Nhà}

\Cref{tab:mono-kb2} trình bày số liệu steady-state trong 30 phút.

\begin{table}[H]
\centering
\caption{Số liệu MONO-KB2 (steady 30 phút)}
\label{tab:mono-kb2}
\small
\begin{tabular}{p{0.42\textwidth} c p{0.30\textwidth}}
\toprule
\textbf{Metric} & \textbf{Giá trị} & \textbf{Ghi chú} \\
\midrule
Complete Latency (mean Grafana) & 1,13 s & Bao gồm drain phase \\
Complete Latency (giai đoạn ổn định) & 2,75 s & Giai đoạn ổn định \\
Complete Latency (max) & 5,09 s & Cuối session khi publisher dừng \\
Split Bolt Capacity (max) & \textbf{99--189\%} & Bottleneck rõ ràng \\
Forecast-1 Execute Latency (max) & 47,8 ms & Bolt nặng nhất ngoài split \\
Spout Emitted Rate (max) & 296 ops/s & Bị throttle từ 400 input \\
DB Rows Written & 26.965 & Trong 30 phút \\
End-to-End Latency p50 & 655 s & Measurement artifact \\
End-to-End Latency p95 & 1.258 s & Measurement artifact \\
\bottomrule
\end{tabular}
\end{table}

\textbf{Phân tích.} Số liệu Bảng~\ref{tab:mono-kb2} khẳng định và mở rộng kết luận
về điểm bão hoà từ \Cref{tab:mono-kb1}. Ba giá trị complete latency cho thấy bản
chất của tình trạng quá tải: giá trị trung bình do Grafana báo (1,13 s) bị
\textit{kéo xuống} bởi giai đoạn drain cuối phiên, trong khi giá trị ở giai đoạn ổn
định thực sự (2,75 s) và đỉnh (5,09 s) mới phản ánh đúng độ trễ mà người dùng cảm
nhận khi hệ thống chịu tải 400 msg/s liên tục. Mấu chốt nằm ở dòng
\code{Split Bolt Capacity} đạt \textbf{189\%}: bolt phân loại phải làm việc gần gấp
đôi khả năng bền vững, nên backpressure kéo \code{Spout Emitted Rate} xuống chỉ còn
\textbf{296 ops/s} dù đầu vào là 400 msg/s --- tức khoảng một phần tư lượng dữ liệu
bị dồn ứ hoặc trễ. Hai dòng End-to-End Latency (p50 = 655 s, p95 = 1.258 s) tuy
lớn nhưng là \textit{artifact đo lường} của ngữ nghĩa \code{REPLACE INTO}
(\Cref{ch:sosanh}), không nên diễn giải là độ trễ thực. Tổng hợp lại, ở quy mô 40
nhà kiến trúc Monolithic không còn hoạt động ở chế độ ổn định mà liên tục ở trạng
thái quá tải có kiểm soát bằng backpressure --- đây chính là động lực cho kiến trúc
Fog.

Trên biểu đồ topology throughput của Grafana
(\Cref{fig:mono-grafana-kb2}), hệ thống thể hiện hình \textit{bell curve}: thông
lượng tăng dần trong giai đoạn nạp tải, đạt đỉnh, rồi giảm dần trong giai đoạn
\textit{drain} sau khi publisher dừng --- khi Storm tiếp tục xử lý lượng tuple
còn tồn trong hàng đợi nội bộ. Hình dạng này là bằng chứng trực quan cho hiện
tượng dồn ứ: nếu hệ thống xử lý kịp thời gian thực, throughput sẽ phẳng theo tải
đầu vào chứ không kéo dài thành đuôi drain sau khi publisher đã dừng.

\begin{figure}[H]
\centering
\begin{subfigure}{0.49\textwidth}
  \centering\grafana{images/mono/s2_steady_bolt_capacity.png}
  \caption{Bolt capacity steady}
\end{subfigure}\hfill
\begin{subfigure}{0.49\textwidth}
  \centering\grafana{images/mono/s2_steady_topology_throughput.png}
  \caption{Topology throughput (bell curve)}
\end{subfigure}
\caption{MONO-KB2 (Grafana): bolt capacity và topology throughput ở trạng thái ổn định.}
\label{fig:mono-grafana-kb2}
\end{figure}

\section{Tóm tắt điểm mạnh/yếu Monolithic}

\begin{table}[H]
\centering
\caption{Điểm mạnh và hạn chế của kiến trúc Monolithic}
\label{tab:mono-procon}
\small
\begin{tabular}{p{0.46\textwidth} p{0.46\textwidth}}
\toprule
\textbf{Ưu điểm} & \textbf{Hạn chế} \\
\midrule
Kiến trúc đơn giản nhất, dễ triển khai & Bão hoà sớm tại $\sim$20 nhà (200 msg/s) \\
Không cần phần cứng biên & Bolt\_split overload 169\% tại 400 msg/s \\
Tập trung quản lý, dễ giám sát & Complete latency 5,28 s tại h40 \\
Phù hợp quy mô nhỏ ($\le$10 nhà) & Cloud offline $\to$ 100\% mất dữ liệu \\
                                 & Tốn băng thông WAN (raw streaming) \\
\bottomrule
\end{tabular}
\end{table}

\textbf{Nhận định.} \Cref{tab:mono-procon} cho thấy kiến trúc Monolithic không phải
là lựa chọn tồi một cách tuyệt đối: với quy mô nhỏ (dưới 10 nhà), sự đơn giản,
không cần phần cứng biên và dễ giám sát của nó là những ưu điểm thực sự đáng giá.
Vấn đề chỉ xuất hiện khi quy mô tăng: cả ba hạn chế cốt lõi --- bão hoà sớm, mất
dữ liệu khi Cloud offline, và tốn băng thông WAN --- đều là hệ quả trực tiếp của
việc \textit{tập trung toàn bộ xử lý ở một điểm và đẩy mọi bản tin thô lên đó}.
Đây chính xác là ba điểm mà kiến trúc Fog nhắm tới giải quyết, và các chương tiếp
theo sẽ định lượng mức độ cải thiện trên từng phương diện.

%% ############################################################################
\chapter{KẾT QUẢ HỆ THỐNG FOG COMPUTING}
\label{ch:ketqua-fog}
%% ############################################################################

\section{FOG-KB1-A --- Ramp Scalability 8 Gateway}

Môi trường gồm 7 container Docker và 1 Raspberry Pi 3, mỗi gateway phụ trách 5 nhà
theo thứ tự kích hoạt gw-01 $\to$ gw-08. \Cref{tab:fog-kb1a} trình bày số liệu.

\begin{table}[H]
\centering
\caption{Số liệu FOG-KB1-A (8 gateway)}
\label{tab:fog-kb1a}
\scriptsize
\setlength{\tabcolsep}{4pt}
\begin{tabular}{c c c c c c c c c c}
\toprule
\textbf{Nhà} & \textbf{Msg/s} & \textbf{Cloud Cap.} & \textbf{Cloud C.Lat.} & \textbf{GW Exec} & \textbf{GW Flush} & \textbf{CPU/GW} & \textbf{RAM/GW} & \textbf{WAN} & \textbf{Tuples} \\
\midrule
1  & 10  & 0\%            & N/A       & 0,032 ms & 93,8 ms  & 11,8\% & 182 MB & 15,0  & 2.527 \\
5  & 50  & 0\%            & 4,5 ms    & 0,073 ms & 100,3 ms & 16,4\% & 191 MB & 40,2  & 14.659 \\
10 & 100 & \textbf{91,2\%}& 14.520 ms & 0,068 ms & 540,2 ms & 11,4\% & 194 MB & 88,9  & 38.480 \\
20 & 200 & 4,6\%          & 34.974 ms & 0,067 ms & 307,6 ms & 13,1\% & 229 MB & 159,4 & 85.884 \\
40 & 400 & \textbf{0\%}   & 84.499 ms & 0,044 ms & 614,4 ms & 10,6\% & 239 MB & 332,8 & 173.249 \\
\bottomrule
\end{tabular}
\end{table}

\textbf{Phân tích.} Bảng~\ref{tab:fog-kb1a} cho thấy đặc tính nổi bật nhất của
kiến trúc Fog 8 gateway: khi tải tăng gấp 40 lần (từ 10 lên 400 msg/s), tải tính
toán ở \textit{biên} gần như không đổi --- GW Execute Latency luôn dưới
\textbf{0,1 ms}, CPU mỗi gateway dao động quanh 10--16\% và RAM chỉ tăng nhẹ từ
182 lên 239 MB. Điều này khẳng định cơ chế tích luỹ gia tăng (Cải tiến~1) đã
chuyển độ phức tạp xử lý từ ``theo từng bản tin'' sang ``theo từng lần flush'',
khiến chi phí trên mỗi bản tin gần như bằng không. Ở phía Cloud, điểm đáng chú ý
là Cloud Capacity \textit{không} tăng đơn điệu theo tải mà đạt đỉnh 91,2\% tại h10
rồi \textbf{giảm về 0\% tại h40} --- một hành vi phản trực giác sẽ được giải thích
ngay dưới đây. Cột WAN tăng tuyến tính nhưng rất chậm (15 $\to$ 332,8 KB/min),
thấp hơn một bậc độ lớn so với raw streaming của Monolithic. Tóm lại, ở mức tải
cao nhất, nút thắt cổ chai đã được dời khỏi cả biên lẫn Cloud.

\textbf{Lưu ý:} Cloud Complete Latency cao (14--84 s) là artifact của ngữ nghĩa
cumulative trong \code{REPLACE INTO} --- \code{updatedAt} luôn bằng lần ghi cuối,
nên đo từ sensor đến lần ghi cuối. Độ trễ thực của lần ghi đầu tiên xấp xỉ 60 s
(chu kỳ flush).

\Cref{fig:fog-kb1a-cap} thể hiện cloud bolt capacity: spike 91,2\% tại h10 do JVM
cold-start, tự phục hồi về 0\% tại h40.

\begin{figure}[H]
\centering
\begin{tikzpicture}
\begin{axis}[
  width=0.82\textwidth, height=6.2cm,
  ybar, bar width=18pt,
  symbolic x coords={1,5,10,20,40}, xtick=data,
  xlabel={Số nhà}, ylabel={Cloud Bolt Capacity (\%)},
  ymin=0, ymax=100, enlarge x limits=0.15,
  nodes near coords, point meta=explicit symbolic,
  every node near coord/.append style={font=\scriptsize},
]
\addplot[fill=fog8Green!70]
  coordinates {(1,0) [0\%] (5,0) [0\%] (10,91.2) [91,2\%] (20,4.6) [4,6\%] (40,0) [0\%]};
\draw[-{Stealth}, monoRed, thick] (axis cs:20,70) -- (axis cs:10,88)
  node[pos=0, right, font=\scriptsize, text width=3.2cm]{JVM cold-start spike, tự phục hồi về 0\% tại h40};
\end{axis}
\end{tikzpicture}
\caption{FOG-KB1-A: Cloud Bolt Capacity (8 gateway).}
\label{fig:fog-kb1a-cap}
\end{figure}

\Cref{fig:fog-kb1a-exec} cho thấy gateway execute latency phẳng dưới 0,1 ms xuyên
suốt mọi mức tải --- minh chứng cho tích luỹ gia tăng siêu nhẹ tại biên.

\begin{figure}[H]
\centering
\begin{tikzpicture}
\begin{axis}[
  width=0.82\textwidth, height=5.8cm,
  symbolic x coords={1,5,10,20,40}, xtick=data,
  xlabel={Số nhà}, ylabel={GW Execute Latency (ms)},
  ymin=0, ymax=0.1,
  nodes near coords, every node near coord/.append style={font=\scriptsize},
]
\addplot[fog8Green, very thick, mark=*] coordinates {(1,0.032) (5,0.073) (10,0.068) (20,0.067) (40,0.044)};
\end{axis}
\end{tikzpicture}
\caption{FOG-KB1-A: Gateway Bolt Execute Latency --- sub-millisecond toàn bộ tải.}
\label{fig:fog-kb1a-exec}
\end{figure}

\textbf{Staircase activation:} các gateway được kích hoạt tuần tự theo ramp:
gw-01 (h01) $\to$ gw-02 (h10) $\to$ \dots\ $\to$ gw-08 (h40). \textbf{Anomaly JVM
cold-start tại h10:} khi hai gateway flush gần đồng thời cùng kết nối JDBC còn
``lạnh'', cloud bolt spike lên 91,2\% capacity, rồi tự phục hồi nhờ JIT
compilation làm nóng đường xử lý. \Cref{fig:fog-kb1a-grafana} là ảnh Grafana
tương ứng.

\begin{figure}[H]
\centering
\begin{subfigure}{0.49\textwidth}
  \centering\grafana{images/fog/kb1a/cloud_bolt_capacity_600s.png}
  \caption{Cloud bolt capacity}
\end{subfigure}\hfill
\begin{subfigure}{0.49\textwidth}
  \centering\grafana{images/fog/kb1a/all_gateways_bolt_capacity.png}
  \caption{Capacity của 8 gateway}
\end{subfigure}
\caption{FOG-KB1-A (Grafana): cloud bolt capacity và capacity của 8 gateway.}
\label{fig:fog-kb1a-grafana}
\end{figure}

\section{FOG-KB1-B --- Ramp Scalability 1 Gateway (Raspberry Pi 3)}

Môi trường: một Raspberry Pi 3 phụ trách toàn bộ 40 nhà. \Cref{tab:fog-kb1b}
trình bày số liệu.

\begin{table}[H]
\centering
\caption{Số liệu FOG-KB1-B (1 gateway)}
\label{tab:fog-kb1b}
\scriptsize
\setlength{\tabcolsep}{4pt}
\begin{tabular}{c c c c c c c c c c}
\toprule
\textbf{Nhà} & \textbf{Msg/s} & \textbf{Cloud Cap.} & \textbf{Cloud C.Lat.} & \textbf{GW Cap.} & \textbf{GW Exec} & \textbf{GW Flush} & \textbf{CPU} & \textbf{RAM} & \textbf{WAN} \\
\midrule
1  & 10  & \textbf{35,3\%}  & 3,5 ms    & 0,04\% & 0,045 ms & 927,5 ms  & 10,9\% & 186 MB & 14,0 \\
5  & 50  & \textbf{21,2\%}  & 3,5 ms    & 0,14\% & 0,023 ms & 661,2 ms  & 10,2\% & 194 MB & 41,8 \\
10 & 100 & $\sim$0\%        & 8.729 ms  & 0,19\% & 0,049 ms & 967,9 ms  & 13,1\% & 211 MB & 88,8 \\
20 & 200 & $\sim$0\%        & 46.123 ms & 0,80\% & 0,048 ms & 1.061 ms  & 14,9\% & 250 MB & 164,5 \\
40 & 400 & \textbf{318,9\%} & 93.387 ms & 0,95\% & 0,030 ms & 759,2 ms  & 11,2\% & 288 MB & 396,1 \\
\bottomrule
\end{tabular}
\end{table}

Tại h40, một gateway gửi $40 \text{ nhà} \times 5 \text{ windows} = 200$
records mỗi batch vào \code{Bolt\_cloudMerge}; thực thi \code{REPLACE INTO} 200
dòng mất khoảng 47,8 s $\to$ \textbf{318,9\%} capacity. Điểm nghẽn nằm ở Cloud,
không phải biên (GW capacity \textless 1\%). Cột $\sim$0\% tại h10, h20 phản ánh
thời điểm Prometheus lấy mẫu rơi vào giữa hai burst; Grafana xác nhận spike xuất
hiện đột ngột tại h40 chứ không tăng đều.

\Cref{fig:fog-kb1b-grafana} là ảnh Grafana KB1-B.

\begin{figure}[H]
\centering
\begin{subfigure}{0.49\textwidth}
  \centering\grafana{images/fog/kb1b/cloud_bolt_capacity_600s.png}
  \caption{Cloud bolt capacity (spike 318,9\% tại h40)}
\end{subfigure}\hfill
\begin{subfigure}{0.49\textwidth}
  \centering\grafana{images/fog/kb1b/summary_stat_panels.png}
  \caption{Summary stat panels}
\end{subfigure}
\caption{FOG-KB1-B (Grafana): cloud bolt capacity và bảng tổng hợp.}
\label{fig:fog-kb1b-grafana}
\end{figure}

\Cref{fig:kb1a-vs-kb1b} so sánh trực tiếp cloud bolt capacity giữa 8 GW và 1 GW.

\begin{figure}[H]
\centering
\begin{tikzpicture}
\begin{axis}[
  width=0.85\textwidth, height=6.5cm,
  ybar, bar width=12pt,
  symbolic x coords={1,5,10,20,40}, xtick=data,
  xlabel={Số nhà}, ylabel={Cloud Bolt Capacity (\%)},
  ymin=0, ymax=340, enlarge x limits=0.15,
  legend style={at={(0.02,0.98)}, anchor=north west, font=\scriptsize},
]
\addplot[fill=fog8Green!70] coordinates {(1,0) (5,0) (10,91.2) (20,4.6) (40,0)};
\addplot[fill=fog1Orange!80] coordinates {(1,35.3) (5,21.2) (10,0) (20,0) (40,318.9)};
\legend{Fog 8 GW, Fog 1 GW}
\draw[dashed, black] (axis cs:1,100) -- (axis cs:40,100)
  node[pos=0.5, above, font=\scriptsize]{Ngưỡng overload 100\%};
\end{axis}
\end{tikzpicture}
\caption{So sánh Cloud Bolt Capacity: 8 GW tự phục hồi về 0\% tại h40, 1 GW overload 318,9\%.}
\label{fig:kb1a-vs-kb1b}
\end{figure}

\Cref{fig:res-h40} so sánh tài nguyên tại h40 (lấy 8 GW làm mốc 100\%).

\begin{figure}[H]
\centering
\begin{tikzpicture}
\begin{axis}[
  width=0.82\textwidth, height=6.2cm,
  ybar, bar width=16pt,
  symbolic x coords={RAM tổng, CPU tổng, WAN}, xtick=data,
  ylabel={Giá trị tương đối so với 8 GW (\%)},
  ymin=0, ymax=130, enlarge x limits=0.25,
  legend style={at={(0.98,0.98)}, anchor=north east, font=\scriptsize},
  nodes near coords, every node near coord/.append style={font=\tiny},
]
\addplot[fill=fog8Green!70] coordinates {(RAM tổng,100) (CPU tổng,100) (WAN,100)};
\addplot[fill=fog1Orange!80] coordinates {(RAM tổng,15.1) (CPU tổng,12.2) (WAN,119)};
\legend{Fog 8 GW, Fog 1 GW}
\end{axis}
\end{tikzpicture}
\caption{FOG KB1-A vs KB1-B: so sánh tài nguyên tại h40. RAM 1 GW tiết kiệm 84,9\%, CPU tiết kiệm 87,8\%, nhưng WAN cao hơn 19\%.}
\label{fig:res-h40}
\end{figure}

\section{FOG-KB2-A --- Cloud Offline Recovery: 8 Gateway}

Môi trường: 8 GW, tải cố định 40 nhà. Timeline: warmup 15 phút $\to$ outage 5
phút $\to$ drain $\to$ post-recovery 5 phút. \Cref{tab:fog-kb2a} trình bày kết quả.

\begin{table}[H]
\centering
\caption{Số liệu FOG-KB2-A (offline recovery, 8 gateway)}
\label{tab:fog-kb2a}
\small
\begin{tabular}{p{0.45\textwidth} c}
\toprule
\textbf{Chỉ số} & \textbf{Giá trị} \\
\midrule
Thời gian outage & 300 s (5 phút) \\
Max Queue Depth  & \textbf{200 batch} ($= 8 \times 5 \times 5$) \\
Recovery Time    & \textbf{96 s} \\
Data Loss        & \textbf{0,0\%} \\
Drain Rate       & 2,08 batch/s \\
Tuples xử lý (30 phút) & 182.223 \\
\bottomrule
\end{tabular}
\end{table}

Timeline chi tiết: trong outage, queue tăng theo bậc thang
$0 \to 40 \to 80 \to 120 \to 160 \to 200$ batch; sau khi Cloud online,
\code{drainQueue()} xả 200 batch về 0 trong 96 s (2,08 batch/s).
\Cref{fig:kb2-timeline} minh hoạ.

\begin{figure}[H]
\centering
\begin{tikzpicture}
\begin{axis}[
  width=0.9\textwidth, height=6cm,
  xlabel={Thời gian (phút)}, ylabel={Queue Depth (batch)},
  ymin=0, ymax=220, xmin=0, xmax=27,
  xtick={0,15,20,21.6,26.6}, xticklabels={0,15,20,21:36,26:36},
  ytick={0,40,80,120,160,200},
  legend style={at={(0.02,0.98)}, anchor=north west, font=\scriptsize},
]
\addplot[fog8Green, very thick, mark=*] coordinates {
  (0,0) (15,0) (16,40) (17,80) (18,120) (19,160) (20,200) (21.6,0) (26.6,0)
};
\addlegendentry{KB2-A (8 GW)}
\draw[dashed, gray] (axis cs:15,0) -- (axis cs:15,220) node[above, font=\scriptsize, black]{bắt đầu outage};
\draw[dashed, gray] (axis cs:20,0) -- (axis cs:20,220) node[above, font=\scriptsize, black]{Cloud online};
\node[font=\scriptsize, fog8Green] at (axis cs:18,210) {warmup / outage (queue 0$\to$200) / drain 96s};
\end{axis}
\end{tikzpicture}
\caption{Timeline Store-and-Forward KB2-A: warmup $\to$ outage (queue 0$\to$200) $\to$ drain (200$\to$0) $\to$ post-recovery.}
\label{fig:kb2-timeline}
\end{figure}

8 GW drain chậm hơn 1 GW (96 s so với 71 s) do tranh chấp (contention) tại MQTT
broker Cloud khi tám gateway drain đồng loạt. \Cref{fig:kb2a-grafana} là ảnh
Grafana.

\begin{figure}[H]
\centering
\begin{subfigure}{0.49\textwidth}
  \centering\grafana{images/fog/kb2a/store_forward_queue_depth.png}
  \caption{Store-forward queue depth}
\end{subfigure}\hfill
\begin{subfigure}{0.49\textwidth}
  \centering\grafana{images/fog/kb2a/summary_stat_panels_outage_queue200.png}
  \caption{Summary panel khi outage (queue 200)}
\end{subfigure}
\caption{FOG-KB2-A (Grafana): queue depth và bảng tổng hợp lúc outage.}
\label{fig:kb2a-grafana}
\end{figure}

\section{FOG-KB2-B --- Cloud Offline Recovery: 1 Gateway (Raspberry Pi 3)}

\Cref{tab:fog-kb2b} trình bày kết quả của 1 gateway.

\begin{table}[H]
\centering
\caption{Số liệu FOG-KB2-B (offline recovery, 1 gateway)}
\label{tab:fog-kb2b}
\small
\begin{tabular}{p{0.45\textwidth} c}
\toprule
\textbf{Chỉ số} & \textbf{Giá trị} \\
\midrule
Max Queue Depth  & \textbf{25 batch} ($= 1 \times 5 \times 5$) \\
Recovery Time    & \textbf{71 s} \\
Data Loss        & \textbf{0,0\%} \\
Drain Rate       & 0,352 batch/s \\
Tuples xử lý (30 phút) & 175.879 \\
\bottomrule
\end{tabular}
\end{table}

\Cref{fig:kb2-compare} so sánh KB2-A và KB2-B trên ba chỉ số.

\begin{figure}[H]
\centering
\begin{tikzpicture}
\begin{axis}[
  width=0.82\textwidth, height=6.2cm,
  ybar, bar width=16pt,
  symbolic x coords={Queue (batch), Recovery (s), Drain x10 (batch/s)},
  xtick=data, x tick label style={font=\scriptsize, align=center},
  ylabel={Giá trị}, ymin=0, ymax=210, enlarge x limits=0.25,
  legend style={at={(0.98,0.98)}, anchor=north east, font=\scriptsize},
  nodes near coords, every node near coord/.append style={font=\tiny},
]
\addplot[fill=fog8Green!70] coordinates {(Queue (batch),200) (Recovery (s),96) (Drain x10 (batch/s),20.8)};
\addplot[fill=fog1Orange!80] coordinates {(Queue (batch),25) (Recovery (s),71) (Drain x10 (batch/s),3.52)};
\legend{KB2-A (8 GW), KB2-B (1 GW)}
\end{axis}
\end{tikzpicture}
\caption{So sánh KB2-A vs KB2-B: queue depth, recovery time và drain rate ($\times 10$ để dễ nhìn).}
\label{fig:kb2-compare}
\end{figure}

\section{Phân tích cơ chế Store-and-Forward}

Cơ chế Store-and-Forward được mô hình hoá bằng các công thức:
\begin{align}
Q_{max} &= N_{GW} \times W_{windows} \times \left\lfloor \frac{T_{outage}}{T_{flush}} \right\rfloor \\
\text{Drain rate} &= \frac{Q_{max}}{T_{recovery}} \\
\text{Điều kiện hội tụ:}\quad & \frac{Q_{max}}{T_{rec}} > \frac{N_{GW} \times W}{T_{flush}}
\end{align}

\textbf{Kiểm chứng:}
\begin{itemize}
  \item KB2-A: $Q_{max} = 8 \times 5 \times 5 = 200$ batch (khớp đo);
        drain rate $= 200/96 = 2{,}08 > 0{,}67$ batch/s tích luỹ $\checkmark$
  \item KB2-B: $Q_{max} = 1 \times 5 \times 5 = 25$ batch (khớp đo);
        drain rate $= 25/71 = 0{,}352 > 0{,}083$ batch/s tích luỹ $\checkmark$
\end{itemize}

Cả hai kịch bản đều thoả điều kiện hội tụ với biên độ lớn, giải thích vì sao
queue không bao giờ có nguy cơ tràn trong các đợt outage ngắn và recovery nhanh
hơn nhiều so với thời gian outage.

%% ############################################################################
\chapter{SO SÁNH VÀ PHÂN TÍCH}
\label{ch:sosanh}
%% ############################################################################

\section{So sánh Scalability: Monolithic vs Fog}

\Cref{tab:cl-compare} so sánh complete latency tại các nấc tải.

\begin{table}[H]
\centering
\caption{Complete Latency so sánh tại các nấc tải}
\label{tab:cl-compare}
\small
\begin{tabular}{c c c c c}
\toprule
\textbf{Nhà} & \textbf{Msg/s} & \textbf{MONO} & \textbf{FOG-8GW} & \textbf{FOG-1GW} \\
\midrule
1  & 10  & 55 ms     & N/A (sub-ms)  & 3,5 ms \\
5  & 50  & 97 ms     & 4,5 ms        & 3,5 ms \\
10 & 100 & 290 ms    & 14.520 ms*    & 8.729 ms* \\
20 & 200 & 1.111 ms  & 34.974 ms*    & 46.123 ms* \\
40 & 400 & 5.282 ms  & 84.499 ms*    & 93.387 ms* \\
\bottomrule
\end{tabular}
\end{table}

\noindent\textit{* Chú thích: Complete Latency Fog cao do ngữ nghĩa cumulative
của \code{REPLACE INTO} --- đo từ \code{event\_ts} đến lần ghi cuối. Độ trễ thực
ghi lần đầu xấp xỉ 60 s. Hai hệ thống đo các đại lượng khác nhau --- xem
\Cref{sec:why-latency}.}

\subsection{Tại sao Complete Latency Fog \textgreater\ Mono?}
\label{sec:why-latency}

Monolithic tính complete latency từ thời điểm MQTT đến cho tới khi
\code{Bolt\_forecast} ack (vài giây). Fog lại tính từ \code{event\_ts} của sensor
đến cột \code{updatedAt} --- vốn luôn bằng \code{NOW()} của lần \code{REPLACE INTO}
\textit{cuối cùng}. Vì mỗi flush 60 s lại cập nhật \code{updatedAt}, một bản ghi
sinh lúc 14:20 sẽ có \code{updatedAt} = 14:50 $\to$ latency biểu kiến = 30 phút.
Đây là artifact đo lường, không phải độ trễ xử lý thực. \Cref{tab:fair-compare}
đưa ra so sánh công bằng dựa trên các chỉ số không bị artifact.

\begin{table}[H]
\centering
\caption{So sánh công bằng dựa trên chỉ số không bị artifact}
\label{tab:fair-compare}
\small
\begin{tabular}{p{0.26\textwidth} p{0.18\textwidth} p{0.22\textwidth} p{0.20\textwidth}}
\toprule
\textbf{Tiêu chí} & \textbf{Monolithic} & \textbf{Fog (8 GW)} & \textbf{Fog (1 GW)} \\
\midrule
Bolt Capacity tại h40 & \textbf{189\%} & \textbf{0\%} & \textbf{318,9\%} \\
Bottleneck & Bolt\_split (cloud) & JVM cold-start h10 (tự phục hồi) & Bolt\_cloudMerge (không phục hồi) \\
Spout throttle tại h40 & 199 ops/s (vs 400) & Không bị throttle & Bị throttle do overload \\
Khả năng scale thêm & Không (đã saturate) & Có (còn nhiều headroom) & Không (318,9\%) \\
\bottomrule
\end{tabular}
\end{table}

\Cref{fig:cap-h40} là biểu đồ quan trọng nhất của đồ án: so sánh cloud bolt
capacity tại h10, h20, h40 giữa ba hệ thống.

\begin{figure}[H]
\centering
\begin{tikzpicture}
\begin{axis}[
  width=0.86\textwidth, height=6.8cm,
  ybar, bar width=11pt,
  symbolic x coords={10,20,40}, xtick=data,
  xlabel={Số nhà}, ylabel={Cloud Bolt Max Capacity (\%)},
  ymin=0, ymax=340, enlarge x limits=0.35,
  legend style={at={(0.02,0.98)}, anchor=north west, font=\scriptsize},
  nodes near coords, every node near coord/.append style={font=\tiny},
]
\addplot[fill=monoRed!75]   coordinates {(10,3.9) (20,89.8) (40,169)};
\addplot[fill=fog8Green!75] coordinates {(10,91.2) (20,4.6) (40,0)};
\addplot[fill=fog1Orange!85]coordinates {(10,0) (20,0) (40,318.9)};
\legend{MONO, FOG-8GW, FOG-1GW}
\draw[dashed, black] (axis cs:10,100) -- (axis cs:40,100)
  node[pos=0.5, above, font=\scriptsize]{ngưỡng overload 100\%};
\end{axis}
\end{tikzpicture}
\caption{So sánh Cloud Bolt Capacity: Fog 8 GW duy trì 0\% tại h40 trong khi Monolithic và Fog 1 GW đều overload.}
\label{fig:cap-h40}
\end{figure}

\section{So sánh Tài nguyên tại h40}

\Cref{tab:res-h40} so sánh tài nguyên toàn hệ thống tại 400 msg/s.

\begin{table}[H]
\centering
\caption{Tài nguyên toàn hệ thống tại 400 msg/s}
\label{tab:res-h40}
\small
\begin{tabular}{p{0.20\textwidth} c c c p{0.24\textwidth}}
\toprule
\textbf{Tài nguyên} & \textbf{MONO} & \textbf{FOG-8GW} & \textbf{FOG-1GW} & \textbf{Ghi chú} \\
\midrule
Cloud CPU       & $\sim$80\% & \textless 5\% & \textless 5\% & Fog nhẹ hơn ở cloud \\
RAM tổng edge   & 0 MB       & 1.912 MB & 288 MB & Fog cần edge RAM \\
CPU tổng edge   & 0          & $\sim$90\% & $\sim$11\% & Fog cần edge CPU \\
WAN (KB/min)    & $\sim$2.400* & 332,8 & 396,1 & Fog giảm WAN 86\% \\
Containers      & 1 (cloud)  & 8 (edge)+cloud & 1 (Pi 3)+cloud & \\
\bottomrule
\end{tabular}
\end{table}

\noindent\textit{* Ước tính raw: 400 msg/s $\times$ 100 bytes = 40.000 bytes/s = 2.400 KB/min.}

\textbf{Phân tích sự đánh đổi.} \Cref{tab:res-h40} làm rõ rằng lợi ích về khả năng
mở rộng của Fog \textit{không miễn phí} mà là kết quả của một sự dịch chuyển tài
nguyên có chủ đích. Monolithic dồn toàn bộ gánh nặng lên Cloud (CPU $\sim$80\%,
RAM lên tới vài GB) và không tốn tài nguyên biên; Fog làm ngược lại --- Cloud CPU
giảm xuống dưới 5\% nhờ chỉ còn một bolt ghi đè, đổi lại hệ thống cần tài nguyên ở
tầng biên. Điểm mấu chốt là \textit{bản chất} của tài nguyên biên này: nó rẻ và
phân tán (Raspberry Pi, container nhẹ) thay vì đắt và tập trung (EC2 lớn). So sánh
hai cấu hình Fog cho thấy một đánh đổi tinh tế: bản 8 gateway tốn nhiều RAM/CPU
tổng hơn (1.912 MB, $\sim$90\%) nhưng phân tán trên 8 nút nên mỗi nút rất nhẹ và
có khả năng cách ly lỗi; bản 1 gateway tiết kiệm tài nguyên tổng (288 MB, 11\%)
nhưng dồn rủi ro vào một điểm. Về WAN, cả hai bản Fog đều giảm khoảng 83--86\% so
với raw streaming --- một lợi ích chung không phụ thuộc số gateway. \Cref{fig:wan-compare}
và \Cref{fig:ram-breakdown} trực quan hoá hai khía cạnh này.

\begin{figure}[H]
\centering
\begin{tikzpicture}
\begin{axis}[
  width=0.82\textwidth, height=5.4cm,
  xbar, bar width=18pt,
  symbolic y coords={FOG-8GW, FOG-1GW, MONO raw}, ytick=data,
  xlabel={WAN (KB/min)}, xmin=0, xmax=2750,
  xtick={0,500,1000,1500,2000,2500},
  enlarge y limits=0.3,
  yticklabel style={font=\small},
  nodes near coords, every node near coord/.append style={font=\scriptsize},
]
\addplot[fill=accentBlue!70] coordinates {(332.8,FOG-8GW) (396.1,FOG-1GW) (2400,MONO raw)};
\end{axis}
\end{tikzpicture}
\caption{WAN Traffic tại h40: Fog giảm $\sim$86\% so với Monolithic raw.}
\label{fig:wan-compare}
\end{figure}

\begin{figure}[H]
\centering
\begin{tikzpicture}
\begin{axis}[
  width=0.82\textwidth, height=6cm,
  ybar stacked, bar width=26pt,
  symbolic x coords={MONO, FOG-8GW, FOG-1GW}, xtick=data,
  ylabel={RAM (MB)}, ymin=0, ymax=5200, enlarge x limits=0.35,
  legend style={at={(0.98,0.98)}, anchor=north east, font=\scriptsize},
]
\addplot[fill=accentBlue!60] coordinates {(MONO,4096) (FOG-8GW,1024) (FOG-1GW,1024)};
\addplot[fill=fog1Orange!75] coordinates {(MONO,0) (FOG-8GW,1912) (FOG-1GW,288)};
\legend{Cloud RAM, Edge RAM}
\end{axis}
\end{tikzpicture}
\caption{Phân tách RAM: Monolithic dồn lên Cloud (RAM lớn); Fog phân bổ giữa Cloud (nhẹ) và biên.}
\label{fig:ram-breakdown}
\end{figure}

\section{Traffic Reduction}

\Cref{tab:traffic-reduction} tổng hợp mức giảm lưu lượng.

\begin{table}[H]
\centering
\caption{Traffic Reduction so với Monolithic}
\label{tab:traffic-reduction}
\small
\begin{tabular}{p{0.40\textwidth} c c}
\toprule
\textbf{Metric} & \textbf{FOG-KB1-A (8 GW)} & \textbf{FOG-KB1-B (1 GW)} \\
\midrule
Acked Tuple Reduction          & 65,8\%  & \textbf{95\%} \\
Emitted/Transferred Reduction  & 52,5\%  & \textbf{92,5\%} \\
WAN Traffic (vs raw 2.400)     & \textbf{$-86,1\%$} (332,8) & $-83,5\%$ (396,1) \\
\bottomrule
\end{tabular}
\end{table}

Tuy 1 GW có tỷ lệ giảm acked/emitted cao hơn (do chỉ một luồng MQTT duy nhất,
gọn hơn), nhưng tổng WAN của 1 GW lại \textit{cao hơn} 8 GW 19\%: với batch lớn
chứa nhiều khoá khác nhau, tỷ lệ nén GZIP kém hơn so với nhiều batch nhỏ.
\Cref{fig:traffic-reduction} trực quan hoá.

\begin{figure}[H]
\centering
\begin{tikzpicture}
\begin{axis}[
  width=0.82\textwidth, height=6cm,
  ybar, bar width=16pt,
  symbolic x coords={Acked, Emitted, WAN savings}, xtick=data,
  ylabel={Tỷ lệ giảm (\%)}, ymin=0, ymax=100, enlarge x limits=0.3,
  legend style={at={(0.02,0.98)}, anchor=north west, font=\scriptsize},
  nodes near coords, every node near coord/.append style={font=\tiny},
]
\addplot[fill=fog8Green!75]  coordinates {(Acked,65.8) (Emitted,52.5) (WAN savings,86.1)};
\addplot[fill=fog1Orange!85] coordinates {(Acked,95) (Emitted,92.5) (WAN savings,83.5)};
\legend{8 GW, 1 GW}
\end{axis}
\end{tikzpicture}
\caption{Traffic Reduction: 8 GW vs 1 GW theo từng chỉ số.}
\label{fig:traffic-reduction}
\end{figure}

\section{So sánh Fault Tolerance}

\Cref{tab:fault} so sánh khả năng chịu lỗi.

\begin{table}[H]
\centering
\caption{Khả năng chịu lỗi của ba hệ thống}
\label{tab:fault}
\small
\begin{tabular}{p{0.26\textwidth} p{0.18\textwidth} p{0.22\textwidth} p{0.20\textwidth}}
\toprule
\textbf{Tình huống} & \textbf{MONO} & \textbf{FOG-8GW} & \textbf{FOG-1GW} \\
\midrule
Cloud MQTT offline 5 phút & 100\% mất dữ liệu & 0\% loss, 96s recovery & 0\% loss, 71s recovery \\
Edge khi cloud down       & Không có edge     & Tiếp tục 100\%        & Tiếp tục 100\% \\
1 gateway lỗi             & N/A               & 12,5\% nhà ảnh hưởng  & 100\% nhà ảnh hưởng \\
Restart gateway           & N/A               & Queue giữ data, drain sau restart & Queue giữ data \\
\bottomrule
\end{tabular}
\end{table}

\textbf{Phân tích.} \Cref{tab:fault} cho thấy khác biệt về khả năng chịu lỗi giữa
Monolithic và Fog không phải là khác biệt mức độ mà là khác biệt \textit{bản
chất}. Monolithic không có tầng biên nên khi Cloud mất kết nối, không có nơi nào
lưu giữ dữ liệu --- toàn bộ dữ liệu trong thời gian gián đoạn mất vĩnh viễn. Cả
hai bản Fog đều đạt 0\% mất dữ liệu nhờ Store-and-Forward, và vẫn tiếp tục xử lý
tại biên khi Cloud down. Sự khác biệt quan trọng giữa hai bản Fog nằm ở dòng
\textit{1 gateway lỗi}: với 8 gateway, một sự cố chỉ ảnh hưởng 1/8 = \textbf{12,5\%}
số nhà, thể hiện tính \textit{cách ly lỗi} (fault isolation); ngược lại, với 1
gateway, một sự cố làm sập toàn bộ 40 nhà. Đây là lập luận mạnh nhất ủng hộ cấu
hình nhiều gateway trong các triển khai đòi hỏi độ sẵn sàng cao, dù bản 1 gateway
có thời gian phục hồi nhanh hơn (71 s so với 96 s) do không bị tranh chấp broker
khi drain.

\begin{figure}[H]
\centering
\begin{tikzpicture}
\begin{axis}[
  width=0.82\textwidth, height=6cm,
  ybar, bar width=16pt,
  symbolic x coords={Max Queue, Recovery (s), Data Loss (\%)}, xtick=data,
  ylabel={Giá trị}, ymin=0, ymax=210, enlarge x limits=0.3,
  legend style={at={(0.98,0.98)}, anchor=north east, font=\scriptsize},
  nodes near coords, every node near coord/.append style={font=\tiny},
]
\addplot[fill=fog8Green!75]  coordinates {(Max Queue,200) (Recovery (s),96) (Data Loss (\%),0)};
\addplot[fill=fog1Orange!85] coordinates {(Max Queue,25) (Recovery (s),71) (Data Loss (\%),0)};
\legend{KB2-A (8 GW), KB2-B (1 GW)}
\end{axis}
\end{tikzpicture}
\caption{Recovery Metrics KB2: queue depth, recovery time và data loss.}
\label{fig:recovery-metrics}
\end{figure}

\section{Energy Optimizer --- CPU theo trạng thái}

\Cref{fig:energy} so sánh CPU mỗi gateway ở ba trạng thái, làm rõ lợi ích của cơ
chế tự động pause. Khi chỉ 2/8 gateway có nhà active (nấc h10), 6 gateway nhàn rỗi
được đưa về CPU 0\%.

\begin{figure}[H]
\centering
\begin{tikzpicture}
\begin{axis}[
  width=0.8\textwidth, height=5.6cm,
  ybar, bar width=26pt,
  symbolic x coords={Running+traffic, Running idle, Paused}, xtick=data,
  ylabel={CPU mỗi GW (\%)}, ymin=0, ymax=18, enlarge x limits=0.3,
  nodes near coords, every node near coord/.append style={font=\scriptsize},
]
\addplot[fill=fog8Green!70, point meta=explicit symbolic]
  coordinates {(Running+traffic,13) [13\%] (Running idle,3) [3\%] (Paused,0) [0\%]};
\node[font=\scriptsize, align=center] at (axis cs:Paused,5) {cgroups freezer\\ wake \textless 50ms};
\end{axis}
\end{tikzpicture}
\caption{Energy Optimizer: CPU gateway theo trạng thái hoạt động (chỉ áp dụng FOG-8GW).}
\label{fig:energy}
\end{figure}

\section{Ma trận so sánh toàn diện}

\Cref{tab:matrix} là ma trận so sánh cuối cùng giữa ba hệ thống.

\begin{table}[H]
\centering
\caption{Ma trận so sánh toàn diện MONO vs FOG-8GW vs FOG-1GW}
\label{tab:matrix}
\scriptsize
\renewcommand{\arraystretch}{1.25}
\begin{tabular}{p{0.26\textwidth} c p{0.14\textwidth} p{0.16\textwidth} p{0.16\textwidth}}
\toprule
\textbf{Tiêu chí} & \textbf{Trọng số} & \textbf{MONO} & \textbf{FOG-8GW} & \textbf{FOG-1GW} \\
\midrule
Scalability (cap. h40)   & 25\% & Tệ (169\%) & \textbf{Tốt (0\%)} & Tệ (318,9\%) \\
Traffic Reduction WAN    & 20\% & 0\% & $-86,1\%$ & $-83,5\%$ \\
Fault Tolerance          & 20\% & 0\% phòng vệ & \textbf{0\% loss, 96s} & 0\% loss, 71s \\
Edge Resource (RAM+CPU)  & 15\% & \textbf{Không cần} & 1.912MB, $\sim$90\% & 288MB, 11\% \\
Deployment Complexity    & 10\% & \textbf{Đơn giản nhất} & Phức tạp (8 nút) & Trung bình \\
Bottleneck Risk          & 10\% & Bolt\_split & \textbf{Không (h40)} & Bolt\_cloudMerge \\
Energy Optimization      & bonus& Không có & \textbf{Auto-pause (CPU 0\%)} & N/A (1 GW) \\
\bottomrule
\end{tabular}
\renewcommand{\arraystretch}{1.0}
\end{table}

\textbf{Tổng hợp và kết luận so sánh.} Nếu quy đổi mỗi tiêu chí thành thang điểm
0--10 và lấy trung bình theo trọng số trong \Cref{tab:matrix}, ta có điểm tổng hợp
mang tính minh hoạ: \textbf{FOG-8GW $\approx$ 8,3/10}, \textbf{FOG-1GW $\approx$
5,4/10} và \textbf{MONO $\approx$ 3,2/10}. Ba con số này tóm gọn toàn bộ phần so
sánh: kiến trúc Fog 8 gateway vượt trội ở đúng những tiêu chí có trọng số cao nhất
(khả năng mở rộng, chịu lỗi, giảm WAN) và chỉ thua ở các tiêu chí trọng số thấp
hơn (tài nguyên biên, độ phức tạp triển khai). Đáng chú ý, MONO và FOG-1GW tuy
cùng bị overload tại h40 nhưng vì \textit{hai nguyên nhân trái ngược}: MONO nghẽn
ở khâu phân loại đầu vào (\code{Bolt\_split}), còn FOG-1GW nghẽn ở khâu ghi DB
(\code{Bolt\_cloudMerge}) do dồn toàn bộ 40 nhà vào một batch. Điều này cho thấy
việc \textit{phân tán} tải ra nhiều gateway --- chứ không chỉ riêng việc tách tầng
Fog --- mới là yếu tố quyết định khả năng mở rộng. Bảng xếp hạng vì vậy không nên
đọc là ``Fog luôn tốt hơn Monolithic'', mà là ``Fog phân tán đúng cách mới giải
quyết được bài toán quy mô lớn'', dẫn tới các khuyến nghị cụ thể ở \Cref{sec:khuyennghi}.

\section{Ngưỡng và Khuyến nghị}
\label{sec:khuyennghi}

Dựa trên kết quả đo, đồ án đưa ra các khuyến nghị triển khai:
\begin{itemize}
  \item \textbf{Dùng FOG-8GW khi:} $\ge 20$ nhà, yêu cầu cách ly lỗi (fault
        isolation), băng thông WAN hạn chế, và có sẵn nhiều Raspberry Pi.
  \item \textbf{Dùng FOG-1GW khi:} $\le 10$ nhà (cloud capacity \textless 35\%
        an toàn), chỉ có một Pi 3, ưu tiên đơn giản hoá triển khai.
  \item \textbf{Không dùng MONO khi:} \textgreater 10 nhà --- bắt đầu degradation;
        \textgreater 20 nhà --- overload nghiêm trọng.
\end{itemize}

\Cref{fig:roadmap} phác hoạ lộ trình scalability (có ngoại suy cho 80 và 160 nhà).

\begin{figure}[H]
\centering
\begin{tikzpicture}
\begin{axis}[
  width=0.86\textwidth, height=6.8cm,
  xlabel={Số nhà}, ylabel={Cloud Bolt Capacity (\%)},
  xmode=log, log basis x=2, xmin=1, xmax=160,
  ymin=0, ymax=400,
  xtick={1,5,10,20,40,80,160}, xticklabels={1,5,10,20,40,80*,160*},
  legend style={at={(0.02,0.98)}, anchor=north west, font=\scriptsize},
]
\addplot[monoRed, very thick, mark=square*, dashed] coordinates {(1,45.5)(5,11.2)(10,3.9)(20,89.8)(40,169)(80,340)(160,400)};
\addplot[fog8Green, very thick, mark=*] coordinates {(1,0)(5,0)(10,91.2)(20,4.6)(40,0)(80,5)(160,15)};
\addplot[fog1Orange, very thick, mark=triangle*, dashed] coordinates {(1,35.3)(5,21.2)(10,1)(20,1)(40,318.9)(80,400)(160,400)};
\legend{MONO (ngoại suy), FOG-8GW, FOG-1GW (ngoại suy)}
\draw[dashed, black] (axis cs:1,100) -- (axis cs:160,100) node[pos=0.5, above, font=\scriptsize]{giới hạn 100\%};
\end{axis}
\end{tikzpicture}
\caption{Lộ trình scalability: vùng \textless 100\% là healthy, \textgreater 100\% là overload. (* ngoại suy)}
\label{fig:roadmap}
\end{figure}

%% ############################################################################
\chapter{KẾT LUẬN VÀ HƯỚNG PHÁT TRIỂN}
\label{ch:ketluan}
%% ############################################################################

\section{Kết quả đạt được}

Đồ án đã đạt được năm kết quả chính:

\begin{enumerate}
  \item Thiết kế và triển khai thành công hệ thống Fog Computing trên Apache
        Storm với \textbf{mười cải tiến kỹ thuật} so với kiến trúc Monolithic.
  \item Chứng minh Fog-8GW đạt \textbf{0\%} cloud bolt capacity tại h40 so với
        Monolithic \textbf{169\%} --- loại bỏ hoàn toàn điểm nghẽn trung tâm.
  \item Cơ chế Store-and-Forward đạt \textbf{0\% mất dữ liệu} trong 5 phút outage,
        phục hồi trong \textbf{71--96 s}.
  \item Giảm lưu lượng WAN khoảng \textbf{86\%} (332,8 so với $\sim$2.400 KB/min
        ước tính raw).
  \item Gateway Raspberry Pi 3 xử lý 50 msg/s với execute latency \textbf{\textless 0,1 ms}
        và bolt capacity \textbf{\textless 1\%} --- còn headroom rất lớn.
\end{enumerate}

\section{Hạn chế và điểm chưa hoàn thiện}

\begin{enumerate}
  \item Fog-1GW bị overload tại h40 (318,9\%) --- cần cải thiện
        \code{Bolt\_cloudMerge} bằng batch transaction.
  \item End-to-end latency là artifact đo lường --- cần cơ chế \code{event\_ts}
        riêng cho lần ghi đầu tiên.
  \item Chưa triển khai đủ 8 Raspberry Pi 3 vật lý (mới dùng 1 Pi + 7 Docker).
  \item Chưa đánh giá tiêu thụ điện năng phần cứng của Raspberry Pi 3.
\end{enumerate}

\section{Hướng phát triển}

\begin{enumerate}
  \item Tối ưu \code{Bolt\_cloudMerge}: ghi nhiều dòng trong một transaction để
        giảm overhead trên mỗi dòng.
  \item Thêm cửa sổ 5 phút và 1 giờ để so sánh trực tiếp với 8 cửa sổ của
        Monolithic.
  \item Auto-scaling gateway: tự động tăng số gateway khi tải tăng.
  \item Triển khai hoàn toàn trên cụm Raspberry Pi 3 (Kubernetes k3s).
  \item Mã hoá payload MQTT (TLS) kết hợp GZIP trong môi trường production.
  \item So sánh thêm với Apache Flink và Apache Spark Streaming.
\end{enumerate}

%% ############################################################################
%%  TÀI LIỆU THAM KHẢO
%% ############################################################################
\begin{thebibliography}{99}
\addcontentsline{toc}{chapter}{Tài liệu tham khảo}

\bibitem{bonomi2012} F. Bonomi, R. Milito, J. Zhu, and S. Addepalli,
``Fog computing and its role in the internet of things,'' in
\textit{Proc. 1st MCC Workshop on Mobile Cloud Computing}, 2012, pp. 13--16.

\bibitem{nist2018} M. Iorga \textit{et al.}, ``Fog Computing Conceptual Model,''
NIST Special Publication 500-325, National Institute of Standards and Technology,
2018.

\bibitem{openfog2017} OpenFog Consortium, ``OpenFog Reference Architecture for
Fog Computing,'' OPFRA001.020817, 2017.

\bibitem{shi2016} W. Shi, J. Cao, Q. Zhang, Y. Li, and L. Xu,
``Edge computing: Vision and challenges,'' \textit{IEEE Internet of Things
Journal}, vol. 3, no. 5, pp. 637--646, 2016.

\bibitem{yi2015} S. Yi, C. Li, and Q. Li, ``A survey of fog computing: Concepts,
applications and issues,'' in \textit{Proc. Workshop on Mobile Big Data}, 2015,
pp. 37--42.

\bibitem{storm} Apache Software Foundation, ``Apache Storm Documentation,''
\url{https://storm.apache.org/}, truy cập 2026.

\bibitem{toshniwal2014} A. Toshniwal \textit{et al.}, ``Storm@twitter,'' in
\textit{Proc. ACM SIGMOD}, 2014, pp. 147--156.

\bibitem{mqtt} OASIS, ``MQTT Version 5.0 Specification,'' OASIS Standard, 2019.

\bibitem{psaras2018} I. Psaras \textit{et al.}, ``Mobile data repositories at the
edge,'' in \textit{USENIX HotEdge}, 2018.

\bibitem{refit} D. Murray, L. Stankovic, and V. Stankovic, ``An electrical load
measurements dataset of United Kingdom households from a two-year longitudinal
study,'' \textit{Scientific Data}, vol. 4, 160122, 2017.

\bibitem{prometheus} Prometheus Authors, ``Prometheus: Monitoring system and
time series database,'' \url{https://prometheus.io/docs/}, truy cập 2026.

\bibitem{grafana} Grafana Labs, ``Grafana Documentation,''
\url{https://grafana.com/docs/}, truy cập 2026.

\bibitem{docker} Docker Inc., ``Docker Documentation,''
\url{https://docs.docker.com/}, truy cập 2026.

\bibitem{mysql} Oracle, ``MySQL 8.4 Reference Manual,''
\url{https://dev.mysql.com/doc/}, truy cập 2026.

\bibitem{atzori2010} L. Atzori, A. Iera, and G. Morabito, ``The internet of
things: A survey,'' \textit{Computer Networks}, vol. 54, no. 15,
pp. 2787--2805, 2010.

\bibitem{gubbi2013} J. Gubbi, R. Buyya, S. Marusic, and M. Palaniswami,
``Internet of Things (IoT): A vision, architectural elements, and future
directions,'' \textit{Future Generation Computer Systems}, vol. 29, no. 7,
pp. 1645--1660, 2013.

\bibitem{satyanarayanan2017} M. Satyanarayanan, ``The emergence of edge
computing,'' \textit{Computer}, vol. 50, no. 1, pp. 30--39, 2017.

\bibitem{dastjerdi2016} A. V. Dastjerdi and R. Buyya, ``Fog computing: Helping
the internet of things realize its potential,'' \textit{Computer}, vol. 49,
no. 8, pp. 112--116, 2016.

\bibitem{mahmud2018} R. Mahmud, R. Kotagiri, and R. Buyya, ``Fog computing: A
taxonomy, survey and future directions,'' in \textit{Internet of Everything},
Springer, 2018, pp. 103--130.

\bibitem{naha2018} R. K. Naha \textit{et al.}, ``Fog computing: Survey of trends,
architectures, requirements, and research directions,'' \textit{IEEE Access},
vol. 6, pp. 47980--48009, 2018.

\bibitem{hu2017} P. Hu, S. Dhelim, H. Ning, and T. Qiu, ``Survey on fog
computing: architecture, key technologies, applications and open issues,''
\textit{Journal of Network and Computer Applications}, vol. 98, pp. 27--42, 2017.

\bibitem{yousefpour2019} A. Yousefpour \textit{et al.}, ``All one needs to know
about fog computing and related edge computing paradigms: A complete survey,''
\textit{Journal of Systems Architecture}, vol. 98, pp. 289--330, 2019.

\end{thebibliography}

%% ############################################################################
%%  PHỤ LỤC
%% ############################################################################
\appendix
\addcontentsline{toc}{chapter}{Phụ lục}

\chapter{Giải thích artifact End-to-End Latency}
\label{app:artifact}

Cột \code{updatedAt} trong bảng \code{fog\_device\_data} được gán bằng
\code{NOW()} tại mỗi lệnh \code{REPLACE INTO}. Do gateway gửi \textit{cumulative
totals} mỗi 60 giây, cùng một khoá bản ghi sẽ bị ghi đè nhiều lần, và
\code{updatedAt} luôn phản ánh lần ghi \textit{cuối cùng}. Khi tính End-to-End
latency $= updatedAt - event\_ts$, kết quả là khoảng cách từ thời điểm sự kiện
gốc đến lần cập nhật cuối --- có thể lên tới hàng trăm giây (p50 = 655 s,
p95 = 1.258 s đối với Monolithic; p50 = 408.626 ms đối với Fog-1GW). Đây
\textbf{không} phải độ trễ xử lý thực. Độ trễ ghi lần đầu thực tế xấp xỉ 60--65 s
(chu kỳ flush + xử lý Cloud). Để đo chính xác cần lưu thêm cột \code{firstWrittenAt}
chỉ gán một lần.

\chapter{Code snippets quan trọng}
\label{app:code}

\begin{lstlisting}[language=Java,caption={drainQueue() --- chuyển tiếp hàng đợi đĩa}]
void drainQueue() {
  List<byte[]> batches = readAll("queue.jsonl");
  for (byte[] b : batches) {
    try {
      mqtt.publish("fog-data", b);  // gui lai theo thu tu
      removeEntry(b);               // xoa khi thanh cong
    } catch (MqttException e) {
      break;                        // dung neu cloud van offline
    }
  }
}
\end{lstlisting}

\begin{lstlisting}[language=Java,caption={Bolt\_cloudMerge.execute() --- ghi idempotent}]
public void execute(Tuple t) {
  byte[] gz = (byte[]) t.getValueByField("batch");
  List<Record> recs = parse(gunzip(gz));
  for (Record r : recs) {
    stmt.execute(
      "REPLACE INTO fog_device_data " +
      "(house_id,household_id,device_id,slice_gap,slice_index,value,count,updatedAt) " +
      "VALUES (" + r.toSql() + ", NOW())");
  }
  collector.ack(t);
}
\end{lstlisting}

\begin{lstlisting}[language=Java,caption={Phát hiện bất thường EWMA z-score}]
double mu = ewmaMu.getOrDefault(dev, value);
double var = ewmaVar.getOrDefault(dev, 1.0);
double z = (value - mu) / Math.sqrt(var);
if (Math.abs(z) > 3.0) sendSlackAlert(house, dev, value, z);
// cap nhat EWMA (alpha = 0.1)
ewmaMu.put(dev, 0.9 * mu + 0.1 * value);
ewmaVar.put(dev, 0.9 * var + 0.1 * (value - mu) * (value - mu));
\end{lstlisting}

\chapter{Cấu hình Docker Compose}
\label{app:compose}

\Cref{tab:compose} tóm tắt các service và cổng chính trong
\code{docker-compose.yml}.

\begin{table}[H]
\centering
\caption{Các service và cổng trong Docker Compose}
\label{tab:compose}
\small
\begin{tabular}{p{0.26\textwidth} p{0.16\textwidth} p{0.44\textwidth}}
\toprule
\textbf{Service} & \textbf{Cổng} & \textbf{Vai trò} \\
\midrule
mosquitto      & 1883       & MQTT broker (iot-data, fog-data) \\
fog-gateway    & 8080--8087 & Storm LocalCluster + Prometheus exporter mỗi GW \\
storm-cloud    & 6700+      & Cụm Storm Cloud (Bolt\_cloudMerge) \\
mysql          & 3306       & Lưu trữ fog\_device\_data \\
prometheus     & 9090       & Thu thập metric \\
grafana        & 3000       & Trực quan hoá \\
web-dashboard  & 9000       & Giao diện người dùng \\
\bottomrule
\end{tabular}
\end{table}

\chapter{Prometheus Queries}
\label{app:promql}

Các truy vấn PromQL chính dùng để đo metric:

\begin{lstlisting}[caption={Các PromQL query chính}]
# Cloud bolt capacity
storm_bolt_capacity{topology="fog-cloud", component="Bolt_cloudMerge"}

# Gateway execute latency (EMA)
rate(fog_gateway_execute_latency_sum[1m])
  / rate(fog_gateway_execute_latency_count[1m])

# WAN traffic (KB/min)
rate(fog_gateway_publish_size_bytes[1m]) * 60 / 1024

# Store-and-Forward queue depth
fog_gateway_store_queue_size
\end{lstlisting}

\chapter{Danh sách ảnh Grafana}
\label{app:images}

\Cref{tab:img-map} ánh xạ tên tệp ảnh $\to$ metric $\to$ kịch bản đo.

{\scriptsize
\begin{longtable}{p{0.52\textwidth} p{0.20\textwidth} p{0.14\textwidth}}
\caption{Bản đồ ảnh Grafana đã sử dụng}
\label{tab:img-map}\\
\toprule
\textbf{Tệp ảnh} & \textbf{Metric} & \textbf{Kịch bản} \\
\midrule
\endfirsthead
\toprule
\textbf{Tệp ảnh} & \textbf{Metric} & \textbf{Kịch bản} \\
\midrule
\endhead
\code{mono/s1\_ramp\_bolt\_capacity.png}      & Bolt capacity      & MONO-KB1 \\
\code{mono/s1\_ramp\_complete\_latency.png}   & Complete latency   & MONO-KB1 \\
\code{mono/s2\_steady\_bolt\_capacity.png}    & Bolt capacity      & MONO-KB2 \\
\code{mono/s2\_steady\_topology\_throughput.png} & Topology throughput & MONO-KB2 \\
\code{fog/kb1a/cloud\_bolt\_capacity\_600s.png}  & Cloud capacity  & FOG-KB1-A \\
\code{fog/kb1a/all\_gateways\_bolt\_capacity.png}& GW capacity     & FOG-KB1-A \\
\code{fog/kb1b/cloud\_bolt\_capacity\_600s.png}  & Cloud capacity  & FOG-KB1-B \\
\code{fog/kb1b/summary\_stat\_panels.png}        & Summary         & FOG-KB1-B \\
\code{fog/kb2a/store\_forward\_queue\_depth.png} & Queue depth     & FOG-KB2-A \\
\code{fog/kb2a/summary\_stat\_panels\_outage\_queue200.png} & Summary outage & FOG-KB2-A \\
\code{fog/kb2b/store\_forward\_queue\_depth.png} & Queue depth     & FOG-KB2-B \\
\code{fog/kb2b/summary\_stat\_panels\_outage\_queue25.png} & Summary outage & FOG-KB2-B \\
\bottomrule
\end{longtable}
}

\end{document}
